% =====================================================================
%  Harald Høffding, Filosofien i Sverig
%  In: Nordisk tidskrift för vetenskap, konst och industri, utgifven af Letterstedtska föreningen, [Ny serie] Andra årgången (Stockholm 1879), pp. 30-73.
%  Antiqua.  Danish in the journal's Nordic spelling (å, small nouns, ö beside ø); quotations in Swedish, German, French, Latin and Greek.
%
%  TRANSCRIPTION -- GENERATED; DO NOT EDIT BY HAND.  Høffding's three pieces of 1879-82 (Filosofien i Sverig; Om
%  Pessimismen; the review of Martensen's Jacob Böhme) are built by ONE pipeline (ebook-method/hoeffding-1879-82/) into a
%  combined working file, which split.py cuts into three transcriptions.
%  Rebuild from the repo root:   sh ebook-method/hoeffding-1879-82/rebuild.sh
%
%  STATUS: COMPLETE (2026-10-15): the volume's title page and the whole text, pp. 30-73, with its 19 notes.
%
%  EDITION.  The only printing.  Text of record: Nasjonalbiblioteket's scan f94504a278617acf87fa0f21fb1dcd03 (URN:NBN:no-nb_digitidsskrift_2015011581044_001,
%  ~400 ppi colour JPEG); extract ~/bibliotek/Høffding, Harald/1879-nordisk-tidskrift-2-pp30-73.pdf (SCANS.tsv, sha256), made by extract.py
%  (title, Innehållsförteckning, the article).  Project Runeberg's nordtidskr/1879 shows the same images.  Page range
%  from the Innehållsförteckning („Höffding, H.: Filosofien i Sverig 30“) and the scan (p. 74 opens the Literaturöfversikt).
%
%  ACCURACY (2026-10-15; one collation for all three texts, 68 printed pages, ~28,100 words).
%  (a) No second copy of any of the three was found (Internet Archive, Nasjonalbiblioteket, Project Runeberg, Google Books;
%      HathiTrust not reachable).  The DRAFT is our own tesseract reading (tessdata_best, 300 dpi; dan+swe for Nordisk
%      tidskrift, dan for the others), chosen by measurement (words unattested in the corpus, and for Nordisk tidskrift
%      a majority of three readings: measure.py), laid over the working PDF and turned into LaTeX by draft.py.  The
%      WITNESS, collated word by word (collate/check.py): NB's ALTO (ABBYY) for Nordisk tidskrift and Vor Ungdom, the
%      Archive's tesseract for Ude og Hjemme.  1,213 disagreements in all.
%  (b) Rules (rules.py: ALTO's known faults -- it drops the umlaut and reads ö as b/d, æ as e -- glyph faults against
%      attested words; umlaut-R/R2: ö/ä where ALTO, the draft and Project Runeberg's OCR of the SAME images agree).
%      Runeberg's votes on other differences were WITHDRAWN after the audit (5 of 30 audited rule/vote decisions wrong,
%      all Runeberg votes: Runeberg shares the other readings' faults) and those 183 places read blind at the image.
%      The remaining rule classes: 0 contradicted in the audit.
%  (c) Every other disagreement read at the image BLIND (rounds b1/b2: 746 crops on 51 sheets, 12 readers, 112 of 112
%      planted controls right), plus the spots of three sweeps on the draft (words unattested in the corpus, specks,
%      and ö/ä where Runeberg's OCR has them and draft and witness agree on o/ø/a/å).  Italics read blind (i1: 6
%      sheets).  OPEN 0.
%  (d) THE PRINT'S ORTHOGRAPHY is the journal's Nordic spelling: å for aa, small nouns, and Danish ö BESIDE ø (göre,
%      större, höjere, Boström / forsøgt, løbe, første) -- both OCRs confuse the two; every ö/ø the readings differ on
%      was settled at the image or by the umlaut rules above.  Quotation marks Swedish-style, ” at both ends.
%  (e) Every FOOTNOTE read whole at the image (reader N1; the clipped tops of pp. 71-72 at zooms), set from the reading;
%      every note MARK placed as read (ALTO's stars, the crops) and checked at a zoom (noteplace.py).  STRUCTURE read at the
%      image (reader K): title page, head, section numerals I. and II., the ending (the article is unsigned).
%  (f) LETTERSPACING measured at the image (spacemeasure.py) and every run judged by reader K: two kept (in quotations).
%  (g) Sweeps on the collated text: unattested words, specks, umlauts; every doubtful spot read at a zoom (final-fixes.tsv).
%  PAGE CHECK.  SAMPLED PAGE CHECK (2026-10-15): 4 pages drawn at random from four strata (Filosofien i Sverig pp. 35, 60; Om Pessimismen
%      p. 83; Ude og Hjemme p. 366, left column), 1,595 words transcribed blind at the image by a reader who saw no other
%      text (pagecheck.py): 1 word error (ɔ: read as o:, then mended in all five places of the class), 1 punctuation error
%      („Følelse“. read Følelse,), 2 italic runs (a posteriori, a priori) -- all mended; the reader's 2 other differences were his
%      own misreadings at a zoom.  About 2 errors per 1,000 words before mending.
%  PRINTED AS IS (misprints kept, marked at the line):
%    p. 33: „Tidsrift“ for Tidskrift; blind reader i1
%  CONVENTIONS.  Page turns: % --- p. N --- and \opage{N} (the printed page); a word divided over a page turn is
%  split with %.  Footnotes as \footnote at their marks, numbered per page as printed (*), **), ***)).  \emph{} =
%  letterspacing, \textit{} = italics.  The running heads and the sheet signatures are not transcribed.  Orthography and the ö/ø distinction as printed; the ø rule (CLAUDE.md) applies only where the print shows a plain o.
% =====================================================================
\documentclass[12pt,a4paper,openany]{book}
\usepackage[utf8]{inputenc}
\usepackage[T1]{fontenc}
\usepackage{libertinus}
\usepackage{libertinust1math}
\usepackage[danish]{babel}
\usepackage{textalpha}   % Greek typed directly
\usepackage[protrusion=true,expansion=true]{microtype}
\usepackage{setspace}
\usepackage[margin=1.3in]{geometry}
\usepackage{fancyhdr}
\usepackage{hyperref}
\hypersetup{pdftitle={Filosofien i Sverig}, pdfauthor={Harald Høffding}, hidelinks}
\pagestyle{fancy}
\fancyhf{}
\fancyhead[LE,RO]{\thepage}
\fancyhead[C]{\textit{H. Høffding: Filosofien i Sverig (1879)}}
\fancypagestyle{plain}{\fancyhf{}\fancyhead[LE,RO]{\thepage}}
\setstretch{1.3}
\usepackage{marginnote}
\usepackage{graphicx}
\DeclareUnicodeCharacter{0254}{\reflectbox{c}}   % ɔ (id est)
\makeatletter
\renewcommand{\thefootnote}{\ifcase\value{footnote}\or *\or **\or ***\else\arabic{footnote}\fi)}   % the print marks each page\'s notes *), **), ***)
\makeatother
\newcommand{\opage}[1]{\marginnote{\footnotesize\textit{[#1]}}}
\newcommand{\signature}[1]{\par\smallskip\hfill #1\hspace*{2em}\par}
\pdfstringdefDisableCommands{\renewcommand{\footnote}[2][]{}\renewcommand{\opage}[1]{}}

\begin{document}

\begin{titlepage}
\begin{center}
{\Huge\bfseries NORDISK TIDSKRIFT}\\[0.4em]
{\scriptsize FÖR}\\[0.3em]
{\Large VETENSKAP, KONST {\small OCH} INDUSTRI.}\\[1em]
{\scriptsize UTGIFVEN AF}\\[0.3em]
{\large\bfseries LETTERSTEDTSKA FÖRENINGEN.}\\[0.8em]
{\scriptsize REDIGERAD AF DOCENT CLAES ANNERSTEDT (HUFVUDREDAKTÖR, UPSALA), PROF.\\ C. M. GULDBERG (KRISTIANA) OCH DOCENT JULIUS LANGE (KÖBENHAVN).}\\[0.6em]
\rule{\textwidth}{0.4pt}\\[0.4em]
{\Large 1879.}\\[0.2em]
\rule{\textwidth}{0.4pt}\\
\vfill
% [printer's monogram ornament]
{\large\bfseries STOCKHOLM,}\\[0.3em]
{P. A. NORSTEDT \& SÖNER.}\\[0.6em]
{\small KRISTIANIA.\hfill HELSINGFORS.\hfill KÖBENHAVN.}\\
{\scriptsize ALB. CAMMERMEYER.\hfill G. W. EDLUNDS BOKHANDEL.\hfill G. E. C. GAD.}
\end{center}
\end{titlepage}
\clearpage

% --- p. 30 ---
\setcounter{footnote}{0}%
\opage{30}\begin{center}{\Large\sffamily\bfseries FILOSOFIEN I SVERIG.}\\[0.4em]{\bfseries {\small AF} HARALD HÖFFDING.}\\[0.4em]\rule{3em}{0.4pt}\end{center}\medskip

\begin{small}
1. \textit{Axel Nyblæus:} Den filosofiska forskningen i Sverige från slutet af adertonde
århundradet, framställd i sitt sammanhang med filosofiens allmänna utveckling. Första
delen. I. Lund 1873. --- II. Lund 1875.

2. \textit{Axel Nyblæus:} Trenne filosofiska uppsatser. Ny upplaga. Lund 1878.

3. \textit{Samuel Grubbe:} Filosofiska skrifter i urval. Utgifna af \textit{Axel Nyblæus.} I—II.
Praktiska filosofiens historia. Lund 1876—77.

4. \textit{Christopher Jacob Boström:} Om religionens, vishetens och dygdens begrepp
samt deras förhållande till hvarandra. Öfversättning [fra latin] af \textit{L. Bygdén.}
Upsala 1874.

5. \textit{Waldemar Dons:} Om Boströmianismen. Kristiania 1874.

6. \textit{E---f---t:} Granskning af kandidaten Waldemar Dons' framställning af den
Boströmska filosofien. Upsala 1875.
\end{small}
\begin{center}\rule{3em}{0.4pt}\end{center}

Man har forsøgt at fremstille filosofiens historie som en sammenhængende
kæde af tankeringe, af hvilke den ene griber ind i den
anden. Ja, Hegel mente endog at kunne påvise, at tænkningens historiske
udvikling nöje svarer til begrebernes plads i det logiske
system og fulgte dettes love. Selv om man anser dette for en vilkårlig
hypothese, der kun kan gennemføres ved at göre vold på den
virkelige historie, --- selv om man har indset, at det er meget forskellige
love, der gælde i logiken og i historien, er man dog endnu
tilböjelig til at opfatte filosofiens udvikling således, at den ene filosof
tager fat, hvor den anden har sluppet. Den störste filosof bliver da,
som Stuart Mill har sagt, den, der bedst har studeret sin forgænger.
Det kan nu ikke nægtes, at filosofiens historie frembyder exempler
nok på en sådan jævnt fremskridende udvikling, navnlig indenfor
hver enkelt af de store hovedretninger. Når et indholdsrigt løsen
er givet, når en ny bane er brudt, har tænkningen foreløbig nok at
göre med at løbe linen ud, udtömme de muligheder, som den nye
retning frembyder. Men nu selve de tænkere, der give løsenet og
indlede de nye perioder? De have ingen forgængere haft at studere
og må altså have fået deres impulser fra kilder, der ligge udenfor
den historisk udviklede filosofi. Det er netop det ejendommelige for
geniet, at det ser nye muligheder og ikke lader sig nöje med at
drage slutninger ud fra overleverede principer.

% --- p. 31 ---
\setcounter{footnote}{0}%
\opage{31}Når der altså er steder, hvor den historiske kæde er brudt, nødes
man til at anerkende, at selve filosofien ikke kan være en aldeles
selvstændig, i sig afsluttet bestræbelse. Den har ikke alle sine udviklingsbetingelser
i sig selv, men peger ud over sig selv, navnlig
på de afgörende vendepunkter. Hvilke de ydre vilkår ere, der især
gribe bestemmende ind, har på en måde allerede Plato antydet, når
han (i ”Phaidros”) fremstiller Urania og Kalliope som filosofiens
muser. Det er naturen og historien, der give tænkningen dens indhold
og dens impulser. Langt fra at være autonom, at kunne øse
alt af sig selv, forudsætter filosofien derfor de andre videnskaber,
særlig naturvidenskaben, og den står tillige i den inderligste forbindelse
med menneskelivets ethiske og kulturhistoriske bestræbelser og
opgaver. Vendepunkterne i filosofiens historie ville vise sig betingede
ved, at nye resultater og methoder optages fra naturvidenskaben,
eller at den ethisk-kulturhistoriske stræben giver tænkningen nye
opgaver. Hos Sokrates var det den ethiske interesse, bestræbelsen
for selverkendelse, for klar indsigt i ”det skönne og gode”, der satte
tanken i bevægelse; og at han blev logikens grundlægger, hænger
nöje sammen med, at han blev ethikens grundlægger. Descartes
søgte at indføre den mathematiske methode i filosofien og at udvikle
en mekanisk naturopfattelse. Locke anvendte induktionen og analysen
ved undersøgelsen af den menneskelige erkendelses oprindelse
og udvikling, og han arbejdede på den praktiske filosofis område i
den politiske friheds og tolerancens tjæneste. Kant var fra sin ungdom
stærkt påvirket af Newton, og søgte i dennes ånd at fremstille
erfaringserkendelsen som den eneste gyldige og at påvise dens grundlag
i bevidsthedens oprindelige organisation; og i de resultater, han
nåede, fandt han tillige tilfredsstillelse for det ethiske ideals ret forståede
fordringer.

Men Kalliope og Urania forliges ikke altid vel; de stille ofte
modstridende fordringer og gå hver sin vej, og navnlig har Kalliope
ofte forlangt en privilegeret stilling og vægret sig ved at dele vilkår
med sin søster. Idealismen, der har sin kilde i menneskelivets ethiske
bestræbelser og opgaver og i følelsen af deres værdi, og naturalismen,
der går ud fra den givne virkelighed, som den ved iagttagelse og
analyse søger at gennemtrænge, synes at være uforsonlige modsætninger.
 Idealismen vil måske ikke engang overlade den ydre natur
til sin modstander; den mener selv at kunne göre fyldest også dér
(Plato, Schelling, Hegel), eller den vender sig bort fra al naturforsken
(Sokrates, Fichte). Men i hvert tilfælde forbeholder den sig
det åndelige område, og betragter det som overgreb, når man her vil
anvende lignende methoder som på den ydre naturs område. Denne
modsætning drager sig gennem hele filosofiens historie og er sikkert
den dybeste grund til alle de kampe, hvorom denne beretter. Den
% --- p. 32 ---
\setcounter{footnote}{0}%
\opage{32}ideale stræben, følelsen af livets, særlig det åndelige livs værd ville
og kunne vi ikke opgive. Men på den anden side stiller den iagttagende
og analyserende forsken lige så bestemte krav; den har udstrakt
sine erobringer til alle territorier på den ydre naturs område,
og forlanger nu med stedse större myndighed også at få adgang til
åndens verden --- eller rettere: den har allerede for længst, uden at
spörge om forlov, skaffet sig denne adgang. Den kan i sin utålmodighed
have begået mange overgreb, forhastet anvendt methoder, der
ikke noksom vare afpassede efter det nye områdes ejendommelighed,
og støttet sig til resultater, der ikke havde den sikkerhed og rækkevidde,
man tillagde dem. Men målet, den sigter efter, kan ikke opgives:
den væsenlige enhed og sammenhæng i vor erkendelse, både
hvad methode og resultater angår. Idealismen har ikke mindre gjort
sig skyldig i overgreb, og den har navnlig forvekslet selve de ideale
værdier med de bestemte former, ideer og symboler, som de til visse
tider have iklædt sig. Den vil først høre op med at betragte naturalismen
som sin fjende, når den har givet alle former til pris, så at
de ikke blive skranker for den forskende tanke, og når den har indset,
at det sande, store og höje ikke mister sin glans, fordi det som
alt andet i tilværelsen forudsætter bestemte betingelser og udvikler
sig efter bestemte love.

Kun sjælden have de to tendenser virket hos en filosof med en
sådan harmoni som hos Kant; naturalismen og idealismen gennemtrænge
hos ham hinanden på ejendommelig måde. Dette blev muligt
derved, at han søgte ned til bægge tendensers fælles kilde: den menneskelige
natur. Så længe de få lov at udfolde sig frit hver for sig,
stiger modsætningen; jo mere man nærmer sig udgangspunktet, des
mere konvergere de. Det var netop Kant's geniale opdagelse, at alle
vore tanker og følelser bestemmes og bæres af vor naturs oprindelige
organisation, af vort standpunkt i universet, og at man altså for at
forstå gyldigheden af vor verdensopfattelse, vor erkendelse og vor
vurdering af tingene, må stige ned til dette grundlag og tage det
med i betragtning. Herved bliver det muligt at indstævne de stridende
verdensanskuelser og methoder for en domstol, de alle må
anerkende. I en vis forstand er filosofien derfor først bleven konstitueret
ved Kant. Ingen har som han peget hen på dens inderste
livsnerve.

At de stridende tendenser efter Kant's tid atter have gjort sig
gældende, kan ikke undre, når man ret betænker dels storheden af
den opgave, han stillede, dels den menneskelige utålmodighed, og
dels også, at han selv ikke har udtalt sig fuldstændig klart, og at
han frembød sider, hvorfra han ikke uden ret kunde tages til indtægt
af en enkelt af de stridende parter. Harmonien mellem de to tendenser
er det store mål, hen imod hvilket vi styre; men undervejs
% --- p. 33 ---
\setcounter{footnote}{0}%
\opage{33}vil der være mulighed og brug for mange hypotheser og systemer.
Ethvert dygtigt og originalt forsøg på at gennemføre en verdensanskuelse
fra det ene eller det andet udgangspunkt vil direkte eller
indirekte, positivt eller negativt yde bidrag til fremskridtet.

De \textit{skandinaviske lande} have kun taget liden selvstændig del i
denne store diskussion. Bølgerne fra de store kulturlande have her,
som på andre områder, fortsat sig ind til os og have fremkaldt livlig
interesse og dygtig tilegnelse, men vi have lige så lidt på det filosofiske
som på det religiøse område frembragt nogen genius af første
rang. Dog står Sverig sikkert her over de to broderlande. Intet af
disse kan opvise en sådan række af energiske og selvstændige tænkere,
som Sverig besidder i Thorild, Höijer, Geijer og Boström, for
kun at nævne de betydeligste. Den filosofiske interesse --- som ikke
kunde være hjemløs i et land, der havde frembragt ånder som den
hellige Birgitta og Emanuel Swedenborg --- har endog i Sverig ført
til udvikling af et ejendommeligt system: det Boströmske, som nu er
det herskende på de svenske universiteter, og hvori den filosofiske
bevægelse synes at have fundet sit foreløbige hvilepunkt, ligesom
i sin tid den tyske spekulation fandt sit i Hegel's system.

Den svenske filosofi har kun været lidet kendt udenfor sin hjemstavn.
Hvad særlig Boström angår, da virkede han overvejende ved
det mundlige ord; han udtaler udtrykkelig i et stridsskrift, at hans
opgave er at være ”en levende lærer for den akademiske ungdom,
ikke en skrivende lærer for den svenske almenhed”; til dette sidste
manglede han ”tid, helbred og ævne”, og det var ham nok, at hans
filosofi bedömtes og kritiseredes af hans disciple og tilhørere\footnote[1]{Hr. \textit{Dons} mener, at denne overvejende mundlige overlevering har bevirket en vis autoritetstro hos Boström's tilhængere. Hr \textit{E---f---t} forsvarer det ”levende ord” mod denne beskyldning (pag. 83): ”Tvärtom har det muntliga och lefvande ordet mer lifskraft och är mer verksamt än det döda bokstafsordet. Och detta gäller isynnerhet och i egentlig mening, då ordet har andligt, moraliskt och religiöst, öfver hufvud taget filosofiskt innehåll. Men att det lefvande ordet är mer verksamt än det döda bokstafsordet, måste betyda, att det verkar mer renande och förädlande på sinnet, mer förklarande och upplysande på förståndet, mer stärkande och solliciterande på viljan, för hvilken ett sådant ord är en kraft till lif och det goda. Den verkan, som utgår från ett sådant ord, är sålunda ett högre och klarare sjelfmedvetande”. --- Så rigtigt og smukt dette er, og skönt hr. E---f---t her kan beråbe sig på Plato, så er det dog sikkert, at den mundlige overlevering ikke er gunstig for selvstændig kritik. Tilhøreren vil ikke så let kunne frigöre sig fra lærerens personlighed og det ethiske og æsthetiske indtryk, den gör, medens læseren kan abstrahere fra det personlige forhold og mere uhildet undersøge og bedömme selve sagen.}. Hans % PRINTED AS IS (Tidsrift: for Tidskrift; blind reader i1)
skrifter vare få og lidet udbredte. I Sverig selv er Bostrom's lære
bleven underkastet en lærerig kritik af prof. \textit{Borelius} i Lund, væsenlig
fra et Hegelsk synspunkt. Denne kritik fremkaldte et særdeles
skarpt og hånende svar fra Boström selv. Hvis det gamle ord slår
til: ”Du har uret, Zeus, ti du bliver vred!” --- så må kritiken have
truffet sit mål. --- Det har heller ikke manglet på forsøg fra nabolandene
på at stifte bekendtskab med hans lære. En ung dansk forsk%
% --- p. 34 ---
\setcounter{footnote}{0}%
\opage{34}er, kand. \textit{Kristian Møller,} rejste i året 1869 på Orla Lehmann's
opfordring til Upsala for at studere Boströmianismen; men en altfor
tidlig død hindrede ham i at fremstille resultatet af denne rejse.
Senere har hr. Dons fra Kristiania udgivet en fremstilling af Boströmianismen
og kritiseret denne fra et Hegeliansk synspunkt. Overfor
ham har hr. E---f---t i en smukt skreven afhandling forsvaret den
Boströmske lære. Striden mellem disse to forfattere angår (ligesom
striden mellem Boström og Borelius) væsenlig det spörgsmål: om
Boströmianismen betegner et fremskridt eller en tilbagegang i forhold
til Hegelianismen. Vi ville i det følgende komme til dette spörgsmål
ved undersøgelsen af den svenske spekulations forhold til den
tyske.

Af særlig betydning for almindeligere og grundigere kendskab
til den svenske filosofi er det, at hr. prof. Axel Nyblæus i Lund er
begyndt at skrive dens historie. Prof. Nyblæus, der hører til de
betydeligste forskere af den Boströmske retning, er bekendt som forfatter
til flere retsfilosofiske og religionsfilosofiske afhandlinger. Han
har bl. a. i sin tid givet et særdeles interessant og betydningsfuldt
indlæg i den strid om tro og viden, som for 10—12 år siden førtes
i Danmark. Hans nye værk har de samme fortrin, som allerede udmærkede
hans ældre arbejder: grundighed, klarhed og elegance i
fremstillingen. Af de hidtil udkomne to bind omhandler det første
”det Gustavianska tidehvarfvets män”: Leopold, Rosenstein, Thorild
og Ehrenstein, hvis filosofiske ideer udviklede sig under indflydelse
af Locke og det 18de århundredes franske filosofer. Forfatteren søger
at vise, at allerede disse mænd vidne om, at selvstændige filosofiske
tanker kunne fødes i hans fædreland, og gendrive den mening, at al
filosofisk forsken i Sverig af nogen værdi hovedsagelig har været en
genlyd fra fremmede lande. Det andet bind skildrer Boëthius og
Höijer og deres forhold til Kant og Fichte. Tilbage står skildringen
af Atterbom, Biberg, Grubbe, Geijer og til sidst den hele bevægelses
afslutning i det Boströmske system.

Den fra Locke udgående filosofiske bevægelse havde en afgjort
naturalistisk karakter. Den søgte at gennemføre erfaringsmethoden
på alle områder. Ved sin klarhed og dristighed blev den et mægtigt
våben i hænderne på det 18de århundredes reformatoriske forfattere,
og ligesom dens ophavsmand Locke var knyttet til den engelske revolution
af 1688, således banede den i hine forfatteres hænder vejen
for den franske revolution af 1789. Samtidig med den politiske
reaktion fik den fra Tyskland udgående idealistiske spekulation herredømmet
indenfor filosofien, i det mindste på fastlandet. Det er denne
idealistiske spekulation, som i Sverig udformedes på ejendommelig
vis. Den Boströmske filosofis plads i filosofiens historie er imellem
de fra Kant og Fichte udgående retninger.

% --- p. 35 ---
\setcounter{footnote}{0}%
\opage{35}Skönt prof. Nyblæus's værk endnu ikke er nået frem til Boström,
ville vi dog ved omtalen af de foreliggende to bind ikke kunne
undgå at diskutere den Boströmske filosofi, da forfatteren vurderer
de tidligere filosofiske bevægelser efter denne filosofis principer, ser
den som det mål, der skal nås, og derfor betragter Lockianismens
og Kantianismens historie som naturalismens selvopløsning, dens
successive overgang til idealisme. For den, der har en anden opfattelse
af idealismens og naturalismens indbyrdes forhold, vil den
nyere filosofis historie tage sig anderledes ud. Det vil derfor have
sin interesse at følge forfatteren i hans historiske fremstilling; vi
ville derved komme til at diskutere nogle af de vigtigste principer i
filosofien. Vel giver forf. ikke i de foreliggende bind præmisserne
til sine kritiske vurderinger; de ville først foreligge i sidste bind af
hans værk: fremstillingen af det Boströmske system. Men vi kunne
her foreløbig orientere os ved hjælp af Boström's egne skrifter: foruden
det i overskriften nævnte navnlig hans ”Grundlinier till philosophiska
statslärans propädeutik”.

Et videre bidrag til oplysning om, hvorledes filosofiens historie
tager sig ud fra den svenske spekulations standpunkt, få vi ved det
skrift af Boström's nærmeste forgænger, \textit{Samuel Grubbe,} som prof.
Nyblæus nylig har udgivet. Dets indhold er ethikens historie, men
det giver egenlig med det samme hele filosofiens historie. Fremstillingen
er klar, men noget bred. Forfatteren søger at vise forbindelsen
og vekselvirkningen mellem det ethiskfilosofiske område
og de andre områder for det åndelige og sociale liv. Da værket
skriver sig fra årene 30 og altså ligger forud for Zeller's, Erdmann's
og andres epokegörende arbejder, må det benyttes med varsomhed i
enkelthederne; men det vidner både om pietet hos udgiveren og om
filosofisk interesse hos hans landsmænd, at et sådant foretagende som
udgivelsen af denne for længst afdøde forskers skrifter har kunnet
sættes i værk.

\begin{center}I.\end{center}

Den populære bevidsthed hviler trygt i troen på de ydre tings
absolute virkelighed. Det falder den ikke ind at göre nogen forskel
mellem objekterne for den sanselige iagttagelse og tingene ”i sig
selv”, som de monne være, afset fra menneskelig opfattelse, --- eller
at plage sig med at spörge om, hvorledes tingene ”i sig selv” må
tænkes, eller om de i det hele taget kunne tænkes. Lige så lidt er
der for den nogen vanskelighed ved at antage, at vore erkendelser
opstå derved, at ”tingene” påvirke os. Når vi lære noget, forholde
vi os altså som kar, i hvilke virkelighedens indhold strömmer ind.

% --- p. 36 ---
\setcounter{footnote}{0}%
\opage{36}Det har kostet den filosofiske reflexion hårde kampe at arbejde
sig ud af denne forestillingsmåde, hvilken vi alle ere nødte til at
anvende i praxis, ligesom vi til dagligdags lade solen gå rundt om
jorden, hvor meget vi end ere Kopernikanere i theorien. Og man
kan ikke undre sig over, at den endnu på mange måder gör sig gældende,
selv hos de tænkere, der ivrigst have studeret erkendelsens
udvikling og væsen. Det var John Locke, der først for alvor gjorde
den menneskelige erkendelse til genstand for studium og derved stillede
den nyere filosofis hovedproblemer. Han påviste vel, at erkendelsens
udvikling beror på, at bevidstheden efter ejendommelige love
bearbejder de fra erfaringen hentede enkelte forestillinger; men ved
selve disse enkeltforestillingers optagelse forholdt efter hans mening
bevidstheden sig aldeles passiv. Prof. Nyblæus mener, at man kan
skildre Lockes theori som en bestræbelse for at retfærdiggöre den
populære forestillingsmåde, ”efter hvilken menneskets erkendelse betragtes
som en relation til en ydre genstand, en indvirken på vor
bevidsthed af en udenfor denne liggende virkelighed”. Men hvilken
ret har man til at tilskrive Locke en sådan bestræbelse? At han
ved undersøgelsernes begyndelse deler den populære forestillingsmåde,
kan man ikke med rette bebrejde ham, og man vilde tvært imod
med större ret kunne tillægge ham en bestræbelse for at arbejde sig
ud af den. Hans distinktion mellem primære og sekundære egenskaber
(i hvilken han har Hobbes, Galilei og Descartes til forgængere),
hans halvt skeptiske opfattelse af substansbegrebet og hans fremhæven
af bevidsthedens aktivitet ved dannelsen af de sammensatte forestillinger,
--- alt dette peger bestemt ud over den vulgære realisme.
Når man derfor taler om konsekvenserne af Locke, har man ikke ret
til uden videre at give ham det 18de århundredes materialisme på
hans samvittighed og sige, at først i ”Systéme de la nature” ”lærte
Lockianismens tilhængere at forstå sig selv og indholdet af de principer,
som de hyldede” (Nyblæus, I, 1 p. 60). Med samme ret kunde
man sige om Luther, at han først lærte at forstå sig selv i Thomas
Münzer og Johan fra Leyden. Endnu stærkere udtrykker forf. sig
på et følgende sted: ”De [Locke's franske efterfølgere] drog for lyset,
hvad den Lockeske filosofi gemte i sit skød, og viste, at man ved at
følge dens principer kommer i strid med menneskets hele höjere bevidsthed
og med selve grundvilkårene for dets liv som menneske”
(p. 77). Spörgsmålet er netop, hvilke Locke's principer vare. Han
er en Janus, der peger i to retninger, og forf. siger selv, at det kun
var ved en ukritisk og dogmatisk fremgangsmåde, de franske materialister
nåede til deres resultater.

Grunden til Locke's store indflydelse i det 18de århundrede må
vistnok mindre søges i hans filosofis bestemte resultater end i den
analytiske methode, han først anvendte med energi og i större om%
% --- p. 37 ---
\setcounter{footnote}{0}%
\opage{37}fang. For den stedse stærkere opposition var denne methode, som
allerede Bouillier (”Histoire de la philosophie Cartésienne” II, p. 550)
har bemærket, ”en krigsmaskine mod den gamle sociale bygning”.
Som Locke havde spurgt om vor erkendelses oprindelse og elementer,
således spurgte man nu om de herskende sociale magters berettigelse,
så dem som frugter af bestemte betingelser og søgte måske endog at
eftergöre dem. Man kan ikke forklare den franske materialisme ud
af Locke alene, men må tage de daværende åndelige og sociale forhold
med i betragtning --- et exempel på, hvad vi have fremhævet i
indledningen: at man ej altid kan udlede den ene filosofiske retning
af den anden, men må tage hensyn til faktorer, der ligge udenfor
filosofien. Materialismen selv var et våben, man greb til som det
mest forsvarlige og lettest tilgængelige. Når man læser ”Systéme de
la nature”, får man snarere indtryk af en religiøs og politisk pamflet
end af et filosofisk værk; det har en afgjort polemisk karakter. Det
forekommer mig derfor at være fuldt berettiget, når Alb. Lange i
sin ”Geschichte des Materialismus” forkaster den navnlig ved den
Hegelske skoles indflydelse gængse historiske opfattelse, i følge hvilken
materialismen udledes af erfaringsfilosofien. Den optræder jo
desuden forud for Locke, hos Hobbes, ligesom den i Frankrig optræder
forud for Condillac, hos La Mettrie.

I øvrigt antyder prof. Nyblæus selv, at den rette udvikling af
Locke's principer fører gennem Berkeley og Hume til Kant. Hvad
det gjaldt om, var at få selve begrebet ”erfaring” nærmere undersøgt.
Men i følge prof. Nyblæus er erfaringsmethoden nødvendigvis behæftet
med den populære forestillingsmådes urigtige forudsætninger,
og han vil derfor have denne methode helt ophævet. Historien har
efter hans mening også holdt dom over Locke og hans methode.
Han skildrer gangen i filosofiens historie i det 18de århundrede på
følgende måde.

Berkeley og Hume viste, at vi ikke have nogen ret til at statuere
”ting” udenfor bevidstheden, da vi hele tiden kun have med vore
egne fornemmelser og forestillinger at göre. Erkendelsen kan ikke
uden selvmodsigelse opfattes som en relation til ydre genstande, da
disse netop først blive til for os i og ved erkendelsen.

Men foruden denne negative kritik blev der også nedlagt indsigelse
mod Locke's empirisme fra et standpunkt, der vel var beslægtet
med den, men i sin ethiske tendens hævede sig over den.
Der var en række forskere, som vel ikke formåede videnskabelig at
bemægtige sig principerne for en höjere og idealistisk anskuelse, men
som dog havde anelser om en sådan. Disse forskeres retning betegner
forf. som ”höjere empirisme”, i modsætning til Locke's empirisme
som den lavere. ”Den höjere empiriske realisme har det fælles med
den lavere, at den sætter det absolute i naturen eller i den i rum
% --- p. 38 ---
\setcounter{footnote}{0}%
\opage{38}og tid givne virkelighed. Men medens den lavere empiriske realisme
så godt som udelukkende reflekterer på selve den fysiske realitet
eller materien (stoffet, hvoraf de sanselige ting bestå), så reflekterer
derimod den höjere empiriske realisme på tingenes \textit{formelle} side,
på deres \textit{formbestemthed,} eller med andre ord: på \textit{ordenen, sammenhængen,
lovbundetheden} i verden og principerne for dem” (I, 1 p. 289).
Det mangelfulde også ved den höjere empirisme er, at den ikke søger
den sanselige (relative) verdens grund og endemål i en selvstændig
usanselig virkelighed (I, 1 p. 91). Til denne retning henfører forf.
Rousseau, Herder, Lessing, og i Sverig Thorild og Ehrensvärd. Men för forf. går over til at skildre disse to mænds virksomhed som tænkere,
dvæler han først ved Lockianismens dygtigste repræsentanter i
Sverig: Leopold (f. 1756, d. 1829) og Rosenstein (f. 1752, d. 1824).

Medens den sidstnævnte ganske støtter sig til Locke, Condillac
og Helvetius, er \textit{Leopold} egenlig kun Lockianer i theoretisk henseende.
Fra erfaringen stammer alt erkendelsesindhold, og fornuften
kan ikke gå ud over den. Men vi besidde en anden, höjere sans
foruden den, der er rettet mod det ydre: ligesom den ydre sans
underretter os om fysisk existens, således underretter den indre os
om moralsk fortræffelighed. Den moralsk-æsthetiske følelse er supplement
til den theoretiske metafysik, og først derved bliver egenlig
filosofi mulig, som ikke blot er tjænerinde for erkendelseslysten, men
beskytterinde af menneskehedens höjeste moralske interesser. Vore
höjeste anskuelser kunne ikke göres afhængige af den videnskabelige
forskens resultater. Det er pligtstridigt at betvivle realiteten af de
til den moralske bevidsthed uadskillelig knyttede sandheder (Gud,
forsyn og udødelighed). Til det ”natur-absoluta” (ɔ: ”de yttersta vilkoren
för all möjlig existens”) må altså föjes det ”moral-absoluta”
(ɔ: ”den högsta lagen för all menniskans fria verksamhet”). (I, I p.
134 ff, 142—149).

I Leopold's karakter som tænker finder forf. et udtryk for den
svenske nationalånd, der også indenfor den filosofiske spekulations
område karakteriseres ved en overvejende retning mod det praktiske
og en dermed sammenhængende dyb bevidsthed om personlighedens
og det oversanseliges realitet.

Den harmoni mellem theori og praxis, natur og ånd, som Leopold
ikke kunde nå, var derimod for \textit{Thomas Thorild} (f. 1759, d.
1808) indholdet af den anskuelse, hvori hans revolutionære ånd fandt
hvile. Begejstret af Rousseau og hans natur-evangelium optrådte han
først som den nye poesis (Young's, Klopstock's og Goethe's) forkæmper
mod den af Kellgren repræsenterede franske skole. Utilfreds
med, at Gustav III's regering ikke vilde gå ind på hans forslag om
udvidet trykkefrihed, drog han til England for at grundlægge en
verdensrepublik, der skulde styres af genier og helte. Han udrettede
% --- p. 39 ---
\setcounter{footnote}{0}%
\opage{39}intet for sin hensigt; derimod gjorde han i London bekendtskab med
Swedenborg's skrifter. Efter sin hjemkomst til fædrelandet fortsatte
han sin æsthetiske og politiske polemik, og da han efter Gustav III's
død i et til rigsforstanderen stilet skrift gav sin begejstring for den
franske revolution luft, blev han dömt til landsforvisning. Dog fik
han en årlig understøttelse af regeringen og blev et par år efter
professor i Greifswalde, hvor han levede sine sidste år og forfattede
sine vigtigste filosofiske skrifter (til dels på tysk).

Thorild kaldte selv (i et brev) sin filosofi for naturalisme. Han
ønskede at kunne sige, når han døde: ”Jag såg naturen!” Men
”natur” tages så her ikke i den betydning, i hvilken vi have taget
dette ord i indledningen, når vi talte om naturalisme. Vi forstod ved
”natur” det i erfaringen (den indre og den ydre) givne virkelige, der
underligger bestemte love og udvikler sig under bestemte betingelser.
Thorild derimod anskuer naturen lige så meget digterisk som filosofisk;
hans spekulation er, som Atterbom\footnote[1]{”Svenska siare och skalder”. (Upsala 1847) IV. p. 12.} har udtrykt det: ”en sig
i tanke bekræftende poesi”. Han anskuer naturen som en stor og
dyb harmoni, som en i uendelige, individuelle former åbenbaret helhed.
Virkeligheden er kærlighed i forskellige grader. Alle særskilte
væsener ere kun forskellige grader eller former, i og gennem hvilke
alkærligheden udvikler sit indhold, --- enkelte toner i den store verdensmelodi
(”nur wegen der melodie sterben die tone”, siger han
etsteds). Ethvert enkelt væsen kan kun forstås ved henblik på de
andre væsener, fra hvilke det kun i graden er forskelligt; individualitet
er ”quantum des \textit{quale}”, og det er spekulationens opgave at bestemme”
quantum dati” i hvert enkelt tilfælde, og dermed hvert
enkelt væsens plads i det store hele. Filosofien er ”archimetri”, ɔ:
lære om det evige verdensmål; dens opgave er \textit{”vide quantum!”}
(”känn förhållandet!”).

Den abstrakte tanke kan ikke opfatte tilværelsens levende, artikulerede
system. Derfor polemiserer Thorild mod Kant, Fichte og
Schelling, hvilke han betragter som repræsentanter for den abstrakte
tankeretning, der i stedet for at have blikket vendt mod det hele,
ser så lidet som muligt og altid det ene uden det andet. Der skal
ingen skilsmisse være mellem tanke og fornemmelse; i den sande
viden forenes tankens höjde og fornemmelsens dybde. Thorild fordrer
en slags intellektuel anskuelse, et ”allblick”, der ved et ”götterblick
der allvergleichung” omfatter på én gang det mest universelle
og det mest individuelle, eller rettere alt individuelt. Den, siger han,
der i sin theori om universet ikke går ud fra følelsen, går glip af
kraften, sjælen, guddommen i naturen og ser i denne kun en død
mekanisme. Følelsen, hvor det abstrakte og konkrete endnu ikke
% --- p. 40 ---
\setcounter{footnote}{0}%
\opage{40}have skilt sig fra hinanden, er derfor det første og sidste hos Thorild.
Prof. Nyblæus betegner hans standpunkt klart ved den bemærkning,
at medens følelsen hos Rousseau kun havde psykologisk betydning,
fik den hos Thorild ontologisk betydning. Det er derfor også i höj
grad betegnende, når Geijer siger, at det 18-de århundredes fortörrede
materialisme, der hos franskmændene og særlig hos Holbach vilde göre
livløsheden til existensens almene princip, hos Thorild slog over i
et overgangssystem, der gennem følelsen levendegjorde alt og i denne
alt-levendegörende følelse henrykt øvede en slags naturdyrkelse. Unægteligt
er det, at følelsen egnede sig vel til et sådant overgangssystem.
Da følelsen næmlig, selv når den efter sit væsenlige indhold er rent
åndelig eller usanselig, altid er forenet med et sanseligt element, så
kunne vi forstå, hvorledes den for Thorild kunde vise sig som den
af ham søgte virkelighed, i hvilken kraft og materie, det åndelige og
det legemlige forenedes og sammensmeltedes”. (I, 1 p. 237).

Når Thorild polemiserer mod den idealistiske spekulation, er det
altså ikke for at hævde erfaringsmethoden; han stiller kun følelsens autoritet
i stedet for tankens. Men alligevel finder prof. Nyblæus for megen
empirisme hos ham, navnlig i hans fordring, at tanken skal kontroleres
ved den umiddelbare iagttagelse og ikke har sin vished i sig selv.
Astronomens \textit{tanker} om stjærnehimlen --- siger prof. Nyblæus --- ere
dog noget helt andet end hans \textit{iagttagelser} af den, og disse kunne
ikke udgöre målestokken for bedömmelsen af de første, således som
Thorild fra sit standpunkt måtte fordre det, når han hævder, at de
”rene tanker” kun ere ”rene digte”. (I, 1 p. 217). --- Men er det ikke
netop en grundregel i den induktive logik, at tænkningen, så langt
den end i sine slutninger og beregninger fjærner sig fra det umiddelbar
givne, dog stedse må sammenholdes og sammenlignes med det,
og at erkendelsen af det virkelige først når sin fuldendelse i det lykkelige
møde mellem slutning og iagttagelse? At denne fuldendelse
kun på få punkter er nået, er sandt; men netop astronomien er et
storartet exempel på, hvorledes forskningen ved at støtte sig til nöjagtig
iagttagelse kan nå til en opfattelse, der vender op og ned på
den umiddelbare anskuelse og dog punkt for punkt stemmer med
denne og forklarer den. Vor erkendelse er i det hele og i det enkelte
endnu overvejende hypothetisk, og grunden hertil er netop verifikationens
vanskelighed på så mange punkter. De videnskaber, hvor
subjektive ideer frit kunne udfolde sig, uden straks at standses af
verifikationens strænge politi, ere længst fra deres mål, og der er
ingen tvivl om, at filosofien hører til dem. --- Men spörgsmålet om
forholdet mellem erfaring og spekulation ville vi i en senere sammenhæng
komme tilbage til.

Den af Thorild's filosofiske ideer, som er af störst betydning for
den senere svenske spekulation, er tanken om tilværelsen som et or%
% --- p. 41 ---
\setcounter{footnote}{0}%
\opage{41}ganisk system, en konkret og harmonisk helhed. Denne tanke kan
siges at være den svenske filosofis grundtanke. Men den fremtræder
hos Thorild endnu ikke i den strængt idealistiske form, som den senere
arbejdede sig hen til. ”Havde Thorild”, siger prof. Nyblæus
(I, 1 p. 271), ”mægtet at give en konsekvent videnskabelig udvikling
af denne idé om det hele, vilde han være bleven ført til en rationelt-idealistisk
verdensanskuelse, eller han vilde have set sig nødt til at
erkende, at den sanselige virkelighed ikke svarer til de opstillede
fordringer om en helhed, men forudsætter en anden virkelighed, der
må være rent ideel eller rent åndelig, og derfor også rent fornuftig
og personlig.”

Til en lignende pantheistisk verdensanskuelse, som Thorild's,
førtes \textit{Karl August} Ehrensvärd (f. 1745, d. 1800), hvis betydning særlig
ligger på æsthetikens område. ---

Medens ”den höjere empirisme” kun gennem anelse eller digterisk spekulation kunde vise hen til den ideale verdensanskuelse,
gjordes der et stort positivt fremskridt med Kant's undersøgelser af
den menneskelige erkendelsesævne, som førte til det resultat, at bevidstheden
i sit væsen er aktiv, enhed i mangfoldighed, og at vor erkendelse
derfor vel begynder med sanselig erfaring, men dog ikke
udspringer af den. Herved gjordes der ende på den forudsætning,
hvorfra Locke gik ud: at bevidstheden forholdt sig aldeles passivt ved
erkendelsesstoffets modtagelse.

Men Kant selv er, i følge prof. Nyblæus, endnu til dels hildet i
den populære realisme. Han udleder vore fornemmelser fra ydre
genstandes virken, og antager en ”ding-an-sich” udenfor bevidstheden.
Kun erkendelsens former have deres oprindelige grund i bevidsthedens
væsen; stoffet optages ude fra. Derfor skelner han skarpt mellem
sanseanskuelse og forstand: ved hin alene ere genstande givne, ved
denne tænkes de kun. Her synes da den forudsætning at ligge til
grund, at al virkelighed må være given for sanseanskuelsen.

Vel har Kant, i følge prof. Nyblæus, indlagt sig stor fortjæneste
ved bestemt at skille fornuften fra sansningen og den til sansningen
sig refererende forstand, og han har ved at opfatte fornuften som
bevidstheden om det oversanselige ”udtalt en af de tanker, ved hvilke
han lader det foregående tidsrums empiriske spekulation langt bagved
sig og åbner udsigt til en ny verden”. Men han nægter, at fornuften
ad egne veje kan nå nogen ny erkendelse, og mener, at den kun stiller
et ideal, der fordrer fremskridt ud over det til enhver tid givne og
opnåede. Ligesom Kant, i følge prof. Nyblæus, har uret i at skelne
mellem stof og form, hvad anskuelsesævnen angår, idet alt, der opfattes
i anskuelsesformerne, lige så vel som disse former selv, ligger
præformeret i bevidstheden og kun skal gå over fra mulighed til
virkelighed, fra dunkelhed til klarhed, --- således har han også uret
% --- p. 42 ---
\setcounter{footnote}{0}%
\opage{42}i at opfatte fornuften som en blot formel ævne til at opstille postulater
og idealer, i stedet for at betragte den som en reel ævne, der har det oversanselige
til sit oprindelige indhold. Kant vender det rette forhold
om, ”når han, i stedet for at søge årsagen til den menneskelige fornufts
trang til at gå tilbage til de sidste grunde i menneskets umiddelbare
bevidsthed om ideerne, tvært imod vil grunde ideerne på den
menneskelige erkendelsesævnes indretning og göre dem afhængige
af den. Også dette hysteron-proteron medvirker til Kant's forkastelse
af ideernes konstitutive gyldighed, en forkastelse, der har sin sidste
grund i visse empiriske og realistiske forudsætninger”. (I. 2 p. 35---37).

Som i sin lære om fornuften, således har Kant også i sin Ethik,
sin Religionsfilosofi, sin ”Kritik der Urtheilskraft” givet mange antydninger
af en höjere verdensanskuelse end den, der følger af hans erkendelsestheori.
Men medens den mod realismen vendende side i
hans theori er fuldt udarbejdet, fremtræde disse antydninger kun
sporadisk og få ingen principiel gennemførelse. Kant viser sig ved
bestemmelsen af de vigtigste begreber stedse befængt med empirisme.
Således f. ex. i hans betragtning af organismens begreb, der er så
betydningsfuldt for den ideale verdensanskuelse. Han går her ud fra
den empirisk givne virkelighed, den sanselige organisme, der udvikler
sig i tiden, og bestemmer organismens almindelige begreb efter den.
Men bestemmelsen af organismens karaktermærke som selvudvikling
passer blot på den endelige eller relative, ikke på den fuldkomne eller
absolute organisme. Tænkes derimod organismens enhed at have
fuldkommen magt over hele sit indhold, så kan der ikke blive tale om
nogen udvikling, da al udvikling går ud på at göre bestemmelsernes
potentielle bestemthed ved enheden til aktuel. Udvikling og bevægelse
ligge derfor ikke i organismens begreb, når dette tænkes uden
indskrænkning eller i sin sandhed og fuldkommenhed. Tænkt på
denne måde betyder organismen intet andet end et reelt system eller
en virkelighed, der danner et i sig sluttet, fuldt sammenhængende
hele, d. v. s. en virkelighed, hvis enhed er nærværende i en mangfoldighed
af bestemmelser eller momenter, der omvendt ere nærværende i
enheden og i hverandre. (I, 2 p. 120—121).

Der må derfor, i følge prof. Nyblæus, skelnes mellem en lavere
og höjere side hos Kant. Efter den lavere side deler han Locke's og
empirismens forudsætning om en ydre realitet, efter hvilken erkendelsen
skal rette sig, og af hvilken den påvirkes. Den höjere side består
i den erkendelse, at al virkelighed ikke opgår i den sanselige verden,
men at der gives en anden og fuldkomnere, over rum og tid
hævet virkelighed, end den sanselige, der kun er et fænomen af den.
--- Det er kun den lavere side, der er bleven udviklet af Kants efterfølgere
i Tyskland. Vel opgav Fichte antagelsen af en ”ting-i-sig”
udenfor bevidstheden, men han opfattede dog, ligesom Kant (og em%
% --- p. 43 ---
\setcounter{footnote}{0}%
\opage{43}pirismen), bevidstheden som en tom form, der først skulde producere
sit indhold. Men produktion er bevægelse, udvikling, og forudsætter
altså tiden; følgelig bliver alt åndeligt liv opfattet som endeligt.
Skönt Schelling og Hegel søgte at råde bod på Fichte's ensidighed,
lide deres spekulationer dog af samme væsenlige mangel som hans;
ingen af dem kan opfatte det åndelige liv anderledes end som successivt
udviklende sig, og altså endeligt og relativt. ”Visselig er det
faktisk, at den menneskelige bevidsthed udvikler sit indhold successivt
(i tiden); men dette udtrykker kun en mangel eller begrænsning
hos os og kan følgelig ikke ligge i bevidsthedens egen natur som
sådan eller have nogen gyldighed for den til sit begreb svarende bevidsthed
(den absolute personlighed).” Dette have først svenske forskere
(Grubbe og Boström) indset. (1, 2 p. 122—127).

Förend vi følge forf. videre i hans skildring af, hvorledes Kant's
filosofi blev modtagen og bearbejdet i Sverig, må vi indskyde nogle
bemærkninger i anledning af hans kritik af Kant.

Der er én ting, hvori Kant er aldeles enig med empirismen, næmlig
det, at al vor erkendelse må hentes fra erfaringen. Når tanken vil
gå på egen hånd, bærer den sig ad som duen, der, fordi farten går
rask i den atmosfæriske luft, troede, at den vilde gå endnu bedre i
det tomme rum. Men det er selve begrebet ”erfaring”, Kant underkaster
en prøvelse, som fører til det resultat, at alle vore erfaringer
efter deres form ere bestemte ved vor bevidstheds natur. Ved analyse
af vor erfaringserkendelse søgte han lovene for bevidsthedsfunktionerne,
--- hvilke love ifølge sagens natur må gælde for alt det indhold,
der udvikles og bearbejdes gennem disse funktioner. Medens
Locke kun havde anerkendt bevidsthedens aktivitet ved dannelsen af
de sammensatte forestillinger, påviste Kant en sådan psykologisk aktivitet
ved dannelsen af de mest elementære sansebilleder.

At også kvaliteten af vore fornemmelser, ikke blot de former,
hvori de ordnes, beror på vor organisation, behøvede Kant ikke særlig
at indskærpe, da det var bragt til foreløbig klarhed ved de tidligere
filosofers lære om de sekundære kvaliteter. Når han alligevel skelnede
mellem stof og form, således at stoffet var givet, medens formen
skyldes vor virksomhed, så var det sikkert ikke hans mening, at der
nogensteds i løbet af bevidsthedens udvikling skulde kunne påvises
et aldeles uformet stof. Når bevidstheden --- som han lærte --- i sit
væsen er synthese, sammenfatten til enhed, så må dette karaktermærke
gælde for alt, hvad der kommer frem i bevidstheden. Kant's theori
er her bleven bekræftet ved den nyere sansefysiologi. Lige fra sin
fødselsstund er enhver fornemmelse bragt i forbindelse med andre
fornemmelser og bestemt ved dem. Stoffet i og for sig bliver da
kun noget, hvortil vi nærme os, jo simplere bevidsthedens sammensætning
bliver. Det rene stof er (som en af Kants skarpsindigste
% --- p. 44 ---
\setcounter{footnote}{0}%
\opage{44}disciple, Sal. Maimon, udtrykte det) et differential, der ikke træffes og
ikke kan træffes i nogen virkelig bevidsthed. Enhver virkelig fornemmelse
har altid en vis form. For at der skal være bevidsthed,
må der være flere fornemmelser givne; men de må ikke blot være
som flere; kun ved deres enhed i erindringen er bevidstheden til.
Ethvert bevidsthedsöjeblik sammenfatter stedse flere öjeblikke.

Når Kant alligevel siger, at en genstand kun bliver os given
derved, at den påvirker vor forestillingsævne på en vis måde, lader
han sig af den populære udtryksmåde forlede til at sammenblande et
faktum med dets forklaring. Faktum er det, at vi se, høre etc.; men
at det, vi se og høre, også skulde existere, selv om vi eller andre ikke
se og høre det, derom kunne vi naturligvis umulig vide noget. Men
Kant vilde også uden skade for sin theori have kunnet holde sig til
det simple faktum, at noget er os \textit{givet,} at vi se, høre, i det hele iagttage
det. Dette er det faktum, erkendelseslæren har at gå ud fra og at
analysere. Det virkelige er det, vi opfatte som virkeligt, --- det, vi,
så meget vi end måtte krympe os derved, til sidst må lade stå, som
det er, --- det, vi ikke kunne andet end anerkende. Men det, at vi
”ikke kunne andet”, er åbenbar et rent subjektivt kriterium, og et
andet kan der efter sagens natur ikke gives. For drömmeren er
hans dröm virkelighed; men når han vågner, opdager han, at hans
dröm kun var en illusorisk virkelighed, der blev bestemt ved en mere
omfattende virkelighed og finder sin forklaring i denne. Ingen enkelt
sanseiagttagelse og ingen særlig tankeakt lærer os virkeligheden at
kende; dertil udfordres en bestandig fornyet og aldrig afsluttet sammenfatten
af gamle og nye erfaringer. Så langt vi kunne gå i drömme
uden at støde på skærende modsigelser eller stridende erfaringer, så
langt tro vi også på drömmens virkelighed. Men der kommer et
punkt, hvor tråden brister. Selv den mest systematiske dröm er kun
et fragment i forhold til den totalitet, vor vågne erfaring fremstiller
for os. Og på lignende måde går det med dem af vore ideer, der
ikke have rod i virkeligheden; tidligere eller senere vil deres begrænsning
vise sig, og det opdages, at der er mere i himlen og på
jorden, end vi tænkte i vor filosofi.

Vor erkendelses udvikling er altså en fremskridende overgang
fra dröm til virkelighed, og vort virkelighedsbillede er stedse kun
relativt. Jo flere nye ting der blive os ”givne”, og jo bedre vi forstå
at gennemtrænge og analysere det givne, des bedre erkende vi virkeligheden.
Men gives der da ikke en absolut virkelighed, hvorfra alt
det stammer, der bliver os givet? Her komme vi over på et nyt
område, hvor den metafysiske theori afløser den logisk-psykologiske
analyse; og for at have fast grund under fødderne er det af vigtighed
ikke at sammenblande disse to ting, således som den populære
opfattelse og mange filosofer göre.

% --- p. 45 ---
\setcounter{footnote}{0}%
\opage{45}Det vil altså ses, at man kan bortskære de populært-realistiske forudsætninger
fra Kants theori, uden derfor straks at styrte sig over i
den idealistiske spekulations arme.

Der var også et andet punkt i Kant's filosofi, hvor han syndede
imod sine egne grundsætninger. Ligesom han ved umiddelbar at
statuere en ting i sig gik ud over erfaringen og bevidstheden, således
lod han på ét punkt, næmlig i pligtbuddet, den absolute virkelighed
direkte træde frem i bevidstheden. Men når pligtbuddet er en psykologisk
kendsgerning, må det dog gælde for dets som for alle andre
kendsgerningers vedkommende: at det er et fænomen, der er erfaringens
(her den psykologiske og historiske erfarings) love underlagt. Opfattelse
af en absolut virkelighed er, i følge Kant's egen filosofi, en
\textit{contradictio in adjecto,} da virkeligheden ved at opfattes af os \textit{eo ipso}
er relativ virkelighed. Også her går Kant altså bagved bevidstheden,
--- og det exempel, han havde givet på at bryde den kritiske filosofis
principer, blev trolig fulgt af hans efterfølgere.

Det forekommer mig, at prof. Nyblæus ved sin sondring mellem
en ”lavere” og en ”höjere” side hos Kant er kommen til at overse
selve det centrale hos ham: det kritiske standpunkt, hvori dog hans
egenlige storhed og betydning ligger. At give dybsindige henvisninger
til en idealverden, mystiske fingerpeg til en ”höjere virkelighed”,
er endda ikke så vanskeligt; man har jo hele den Platoniske og Ny-Platoniske
litteratur til sin rådighed, for ikke at tale om theologi og
katekismus. Men til at blive erkendelsestheoriens egenlige grundlægger,
ved logisk-psykologisk analyse at blotte den menneskelige ånds elementære
struktur, dertil hører et geni og udfordres et tankearbejde, som
søger sin lige.

Det er ikke min mening, med visse ivrige Ny-Kantianere at opfatte
den Efter-Kantiske spekulation som et syndefald. I hvert tilfælde
havde Kant selv sat slangen ind i paradiset. Man må for at
være retfærdig mod spekulationens heroer erindre, hvilke misbrug der
kunde drives med Kant's filosofi, aldeles afset fra dens indre mangler. Der
er intet spörgsmål om, at hvad der vandt så mange for Kant's system
og så hurtig skaffede det indgang, nok så meget var hans praktiske som
hans theoretiske resultater. Det var den dybe ethiske alvor, der
fængslede så mange. Men blandt dem var der atter ikke få, som trøstede
sig over umuligheden af enhver spekulativ theologi, over omstyrtningen
af de gode gamle ”beviser for Guds tilværelse”, ved den anvisning,
Kant selv gav på at postulere Gud og udødeligheden af moralske
grunde. De fandt en ny udvej til at redde deres gamle dogmer.
De satte Kant höjest netop dér, hvor han hørte op at være
tænker; ti den Gud, der hos Kant skal bringe dyd og lyksalighed i
harmoni, er et lige så mythisk væsen som ”Demiurgos”, der hos Plato
fremtrylles for at forme materien efter ideerne. Med rette tilråbte
% --- p. 46 ---
\setcounter{footnote}{0}%
\opage{46}den unge Schelling sådanne Kantianere (i sine interessante ”Briefe
über Dogmatismus und Kriticismus”): ”For eder var ikke dogmatismen
men i det höjeste kun den dogmatiske filosofi bleven styrtet!”
--- Det var altså for at hævde filosofiens frihed og selvstændighed
og tilfredsstille fordringen om enhed i verdensanskuelsen, at så mange
begavede ånder slog ind på de spekulative veje. Når man ikke vilde
gå ad erfaringens vej, der ikke fører ud over relativiteten, stod der
kun én vej åben: den idealiserende abstraktion.

Når prof. Nyblæus i sin kritik bebrejder Kant, at han betragter
tidsudviklingen som væsenligt element så vel i bevidsthedens som i
organismens begreb, skönt tiden kun gælder for det endelige, ikke
for det uendelige og absolute, så skylder han først Kant regnskab
for, hvorledes det bliver muligt at göre en sådan overgang fra det
endelige til det uendelige og i det hele at opstille et positivt begreb
om det absolute. Kant vilde til denne kritik ganske simpelt svare,
at han kun har talt om den organisme og den bevidsthed, som vi
\textit{kjende;} først når der kan påvises en bevidsthed uden tidsvilkår, en organisme,
der ikke er udvikling underlagt, vil det være berettiget at
udvide begreberne. Spörgsmålet er, om vi i det hele kunne tænke
det tidløse, --- om ikke forandring, overgang, bevægelse, udvikling
ere fundamentale bestemmelser i vor bevidstheds natur. Og selv om
vi kunde abstrahere fra disse bestemmelser, så er den blotte abstraktion
dog ikke nok til at danne forestillingen om et \textit{virkeligt} væsen.
Ad abstraktionens vej kunne vi nå tilbage til de simple grundforhold
og grundlove, men den kan ikke føre til nogen ny erkendelse af et
virkeligt, som ikke i forvejen er givet. At appellere til ”menneskets
umiddelbare bevidsthed om ideerne” (I, 2 p. 36) hjælper intet, da det
jo ikke er alle, der besidde en sådan bevidsthed, og den i hvert tilfælde
må underkastes en nærmere psykologisk analyse, för man kan
vide, hvad den betyder. Dersom den menneskelige fornuft besad et
sådant oprindeligt indhold og var en ”real ævne”, vilde vi få to slags
logik: en for den höjere og en anden for den lavere erfaring. --- Men
det sidste motiv til at opstille idealet om den fuldendte, over tid og
udvikling hævede virkelighed er --- som det siden vil vise sig --- i
virkeligheden ikke logisk, men ethisk.

Det var også abstraktionens vej, man i Tyskland slog ind på, da
man efter det kritiske mellemspil på ny betrådte de spekulative baner.
Kant's analyse havde ført til opfattelsen af bevidstheden som en oprindelig
synthese. Men synthese forudsætter stedse en flerhed, der
forenes, en dobbelthed, der overvindes. Lad os da --- sagde man ---
søge tilbage til det punkt, hvor dobbeltheden bryder frem, hvor mangfoldigheden
endnu ikke er sat, --- til den rene aktivitet! I al virkelig
bevidsthed finde vi ganske vist kun den relative, bestemte og
modificerede aktivitet; men derfor må spekulationen også bygge på
% --- p. 47 ---
\setcounter{footnote}{0}%
\opage{47}en grund, der ligger ud over bevidstheden. Al virkelig bevidsthed ---
sagde Fichte --- beror på en duplicitet, en sondring mellem subjekt og
objekt. Den rene virksomhed, hvorved først subjekt og objekt blive til,
har derfor selv intet objekt, ingen modstand. For at bevidstheden skal
blive til, erfaring blive mulig, må der sættes skranke og begrænsning;
kun relativ virkelighed er virkelighed for os. Men hvorfor skal da
erfaring være mulig? På dette spörgsmål giver Fichte kun et teleologisk
svar: fordi den uendelige, expansive bevægelse ikke har den
værdi, som bevidsthed og viden. Det at blive sig bevidst er altså et
gode, som skal opnås, og derpå går den evige fornuftproces ud. Under
forskellige former er det det samme svar, Schelling og Hegel også i
sidste instans give. Fichte har så stor interesse, fordi spekulationens
motiver og udgangspunkter hos ham ligge så tydelige.

Vi stå her ved det punkt, hvor vejene skilles. Den vej, der fra
Kantianismen fører ind i spekulationens land: den idealiserende abstraktion,
deler sig i to grene, idet den tyske spekulation dristig tillægger
det absolute en evig selvudvikling, medens den svenske spekulation
betragter det absolute som det evig fuldendte og afsluttede.
Vi ville endnu ikke undersøge, hvilken af disse retninger der har
störst værdi og berettigelse; ti vi have endnu tilbage at betragte den
svenske spekulations historiske udvikling.

Den Kantiske filosofi mødtes i Sverig med skarp kritik fra flere
af de mænd, der havde sluttet sig til Locke's retning; foruden Leopold
optrådte navnlig \textit{Fremling} (d. 1820) med en særdeles skarpsindig
kritik især af Kant's moraltheologi. Han viser, hvorledes Kant her
stedse kommer i modsigelse med sig selv: Postulatet om udødelighed
begrunder Kant ved, at mennesket her i livet aldrig når den
ethiske fuldkommenhed, --- men den kan, som han selv lærer, ikke
nås af noget endeligt væsen, men vedbliver at stå som et ideal. Vel
siger Kant, at tiden ikke har gyldighed for den intelligible verden;
men hvoraf véd han det, når denne intelligible verden ikke kan være
genstand for erkendelse? En personlig udødelighed følger i hvert
tilfælde ikke; ti det kunde jo være, at fuldkommenheden realiseredes
gennem slægtens udvikling, der af Guddommen sammenfattes til enhed.
Og når Kant i sin Ethik vil ophæve egenkærligheden, synes
ønsket om fortsættelse af den personlige tilværelse også at måtte trænges
tilbage. Det er endelig en af Kants methodiske grundsætninger,
at alt i universet skal forklares af naturlige årsager, så at endog de
vildeste fysiske hypotheser ere at foretrække for en appel til det overnaturlige;
men så synes han også ved spörgsmålet om lyksalighedens
forbindelse med sædeligheden at måtte lade sig nöje med en ved naturårsager
betinget moralsk verdensorden.

Andre af Locke's tilhængere blev derimod vundne for den nye
filosofi. Kants første discipel i Sverig var Daniel Boëthius (f. 1751,
% --- p. 48 ---
\setcounter{footnote}{0}%
\opage{48}d. 1810). Det var gennem denne mands lærervirksomhed ved universitetet
i Upsala, at den idealistiske filosofi fik indgang i Sverig og
stødet blev givet til den nationale filosofiske retning. Ikke blot
Höijer, men også Biberg og Grubbe modtog fra ham den første vækkelse
til filosofisk forsken. Han har givet det første udkast til den
ethisk-religiøse verdensanskuelse, der senere fuldendes af de to sidstnævnte
forskere og Boström.\footnote[1]{Af interesse for den almindelige litteraturs historie er følgende bemærkning af prof. Nyblæus (I, 2 p. 236): ”Utan den djupa sedliga anda, som utgick ifrån en Boëthius, en Biberg, en Grubbe och Boström, skulle säkerligen icke ens Geijer kunnat bilda en tillräcklig motvigt emot det upplösande moment, som låg i ”fosforisternas” naturalism.”}

Medens Boëthius er enig med Kant i, at fornuften ikke theoretisk
kan nå det absolute, lægger han mere end Kant vægten på, at
den fordrer harmoni i vort væsen. Denne tanke om harmoni, om det
åndelige liv som en ideel organisme, er jo --- som vi allerede have
set --- den svenske filosofis alfa og omega. Boëthius føres ved den
ud over den Kantiske modsætning mellem theori og praxis, mellem fornuft
og følelse. Han søger at give en genetisk forklaring af det kategoriske
imperativ ved at sc det som form og udtryk for tankens selvvirksomhed,
af hvilken al lov og al harmoni udspringer. Han føres derved til
en sammensmeltning af Kant's ethik og den græske, idet han stiller den
antike beundring for det sædelig skönne (\textit{τὸ καλόν}, \textit{honestum}) sammen
med Kants lære om agtelsen for fornuftloven. Tillige opfattede han det
ethiske som i sin sidste grund ét med det religiøse, sædeligheden som
”et liv i Gud”.

Den egenlige hovedrepræsentant for den Kant-Fichte'ske retning
i Sverig, og i det hele en af de betydeligste svenske filosofer, er \textit{Benjamin}
Höijer (f. 1767, d. 1812). Han var, ligesom Thorild, en ivrig
frihedsmand, og da han havde sluttet sig til en litterær klub i Upsala,
som kaldtes ”Juntan”, og som blev beskyldt for jakobinisme, led han
på grund af ”skuggrädslan hos Gustaf IV Adolf” en række ydmygelser
og tilsidesættelser og blev først ansat som professor efter denne
konges afsættelse, men kort efter faldt han i en sygdom, hvis følger
han ikke forvandt.

”Höijer var”, siger prof. Nyblæus (I, 2 p. 351), ”en af disse overlegne
naturer, der føle sig kaldede til at give love i erkendelsens
verden for den øvrige menneskehed. --- --- Med den overlegne intelligens
og den stærke selvfølelse (som dog overskred tilbørlighedens
grænser) forenede sig hos ham et uforfærdet mod, der ikke betænkte
sig på at drage konsekvenserne af de grundsætninger, som han anså
for rette, og en fast vilje, der ikke trættedes i bestræbelsen for at nå
det mål, han havde sat sig. Betegnende for ham i den første henseende
er et svar, som han skal have givet en velmenende person,
% --- p. 49 ---
\setcounter{footnote}{0}%
\opage{49}der mente at burde henlede hans opmærksomhed på de farer for hans
sjæls salighed, hvormed den frie forsken var forbunden: ”Sök sanningen,
och bure det än til helvetets portar, så klappa på!”

På en rejse i Tyskland havde Höijer gjort Fichte's og Schelling's
bekendtskab, og begge disse tyske filosofer satte ham meget höjt som
tænker. Hans stilling i filosofiens historie er betegnet ved, at han successivt
forudgriber Fichte's, Schelling's og Hegel's ideer og således for sit
eget vedkommende beskriver den samme bane, som i Tyskland med
videre gennemførelse og rigere fylde tilbagelagdes af hine tre tænkere.
Han så nødvendigheden af inderligere sammenhæng og enhed
i filosofien, end Kant, hvem han ellers betragtede som den störste af
alle filosofer, havde bragt til veje; og han vilde, ligesom Fichte, opfatte
bevidstheden ikke blot som grund til erkendelsens form, men
også til dens indhold. Men Fichte bliver stående ved den abstrakte
selvbevidsthed, det rene jeg, og kommer derfor ikke ud over dualismen,
da et jeg er utænkeligt uden i modsætning til et objekt. Derfor
går Höijer tilbage til det princip, der er fælles for subjektiviteten
og objektiviteten, og forudgriber således Schelling's identitetslære.
Medens det for Fichte var umuligt at påvise, hvorledes det rene jeg
kommer til at sætte sig selv skranker og derved blive bevidst, søger
Höijer at påvise begrænsningen som oprindelig given i selve handlingens
begreb. Uden koncentration ingen energi; uden begrænsning
vilde det absolute tabe sig i uendelig expansion. Men han var ikke
enig med Schelling i den stilling, denne anviste naturfilosofien som
sideordnet med transscendentalfilosofien; naturen var for Höijer, ligesom
for Hegel, kun en gennemgangssfære for det absolute. Heller
ikke kunde han slå sig til ro ved Schelling's ”intellektuelle anskuelse”;
han fordrede, ligesom senere Hegel, en begrundelse af denne anskuelse
ved en höjere logik; og det var arbejder i denne retning, der
sysselsatte ham i hans sidste år.\footnote[1]{I sine sidste år betragtede Höijer med stigende utilfredshed den måde, hvorpå de yngre svenske Schellingianere optrådte i deres tidsskrift ”Fosforos”. ”Skönt Höijer's stilling til dem oprindelig blot var undvigende, men langtfra fjendlig, og mere afventende end afvisende, så kunde han, jo mere de røbede, hvad de førte i deres skjold, ikke undgå at frygte, at den slags Schellingianisme, som de søgte at udbrede i vort land, skulde virke ødelæggende for alt virkelig filosofisk studium ved universitetet og bringe filosofien i fuldstændigt vanry. Det var vist nærmest af frygt herfor, at Höijer skrev en afhandling om ”åskådning”, som findes i 2:det bind af hans skrifter, og i hvilken han revser den ”ord- og formelfilosofi, hvortil Schellingianismen var udartet hos mængden af sine tilhængere”. (I, 2 p. 349). --- Endnu på sin dødssæng fralagde han sig med stor iver at have haft nogen del i nævnte tidsskrift.} Endelig kunde han ikke, som Schelling,
betragte kunsten som den höjeste udtryksform for det åndelige
liv; for ham, ligesom for Kant og Fichte, stod den praktisk-ethiske
virksomhed höjest.

Det er en praktisk interesse, som i følge Höijer fører til filosofi,
idet det filosofiske grundproblem er, hvorledes den moralske bevidst%
% --- p. 50 ---
\setcounter{footnote}{0}%
\opage{50}hed og den dermed uadskillelig forbundne frihedsbevidsthed kunne bevares
overfor naturnødvendigheden, eller --- fra en anden side set ---
hvilket forholdet er mellem subjektet og objektet. Derfor er fornuften
ikke tilfredsstillet, för den har nået det absolute, hvor subjekt og objekt
ere ét. Påvisningen af vor erkendelses realitet og forklaringen af
objekternes mulighed falde sammen i den filosofiske konstruktion,
der udleder den relative virkelighed af den absolute og viser, hvorledes
de endelige jeg'er, for hvem objektet eller naturen gör sig gældende,
ikke ere andet end den mangfoldighed af selvindskrænkninger,
hvorunder det absolute åbenbarer sig, og ved hvilke det fuldstændig
bestemmes. --- Den menneskelige viden må være et i sig sammenhængende
hele, et system; og for at kunne være genstand for viden
må tilværelsen selv også være et system. Dette er Höijer's hovedtanke,
ved hvis nærmere gennemførelse han kom til det resultat: at
da den ydre verdens objekter må opfattes som ubevidste produkter af
det absolutes virken, kan man med rette sige, at naturen i sin sidste
grund ikke er andet end den moralske verden set fra en vis side, et
udtryk for menneskets plads indenfor denne verden. Dersom det
menneskelige individ ikke tilhørte en helhed af væsener, der erkende
sædeligheden og skride fremad til den, gaves der ingen moralsk fuldkommengörelse
for det; ti blot indenfor et sådant hele er pligt (ɔ:
”allmän lags giltighet”) mulig. (I, 2 p. 358, 364 ff, 401 f).

Höijer hører endnu ligesom de med ham nærmest beslægtede
tyske tænkere, i følge prof. Nyblæus, til dem, der have udviklet den
”lavere side” i Kantianismen. Han anvender udviklingsbegrebet på
det absolute, lader dette ”komme til sig selv” ved at overvinde
skranker og modsætninger. I virkeligheden er han derfor ikke nået
ud over empirismen og realismen, har ikke nået begrebet om den
sande virkelighed! (I, 2 p. 443).

Des værre miste vi på dette punkt prof. Nyblæus's kyndige vejledning
og må på egen hånd og under egen risiko tilföje nogle træk
til karakteristik af den svenske idealisme i dens fuldendte form. De
vigtigste principer have vi allerede sporadisk mødt i det foregående.

Den spekulative filosofi søger det absolute og finder først hvile i
det. De forskellige spekulative retninger skille sig derfor fra hverandre
ved deres forskellige opfattelse af det absolute. ”Det höjeste
begreb i hele filosofien”, siger prof. Nyblæus (I, 2 p. 424), ”er begrebet
om det absolute eller om den selvstændige og fuldkomne virkelighed,
hvortil den endelige verden på grund af den uselvstændighed
og ufuldkommenhed, hvormed den i en vis henseende er behæftet,
henviser som sin forklaringsgrund”. --- Men den svenske idealisme,
der, forberedt af Biberg, Grubbe og Geijer, fik sin afsluttede form
af Boström (f. 1797, d. 1866), finder afgörende mangler i de fleste
tidligere forsøg på at udvikle dette begreb. Absolut, siger Boström,
% --- p. 51 ---
\setcounter{footnote}{0}%
\opage{51}er det, der er skilt fra og uafhængigt af alt andet, og som i sig selv
indeholder alt det, der udgör dets væsen. Absolut er altså det samme
som fuldkomment; ingen bestemmelse kan tages fra det eller föjes
til det. Absolute væsener kunne derfor ikke have noget udenfor sig,
ikke være underkastede tid og forandring, rum og sted. Selvbevidsthedens
form er den eneste mulige for et absolut væsen; ti i den fuldendte
selvbevidsthed er ét i alt, og alt i ét (\textit{πάντα ἐν παντί}, som
Anaxagoras sagde). (”Om relig.” etc., p. 88—94, 32---34). Som selvstændig
og fuldkommen totalitet besidder det absolute \textit{absolut} uendelighed;
dets fuldendelse er altid virkelig, og det er umuligt at komme
videre. Det sanselige besidder derimod relativ uendelighed, idet dets
fuldendelse aldrig bliver virkelig, men det bestandig er muligt at
komme videre. Tillige besidder det absolute \textit{indre} uendelighed, ɔ:
intet af dets momenter kan opfattes for sig alene, intet kan tænkes,
uden at alle de andre \textit{eo ipso} tænkes; de ere indbyrdes hverandres
momenter. Den absolute fornuft har i sine momenter subjekter og
centra for opfattelsen af sig selv og fornemmer sig selv med hele
sit indhold i den fuldstændige fornemmelse hos et hvilket som
helst af dem. Den har dem således alle nærværende i sig, og den
er selv med hele sit indhold nærværende i ethvert iblandt dem; hvoraf
følger, at de også alle ere nærværende i hverandre. (”Grundlinier”
etc., p. 30). Alle momenter indtage ikke samme rang; der gives lavere
og höjere væsener i det absolute. Men det ene er ikke udenfor det
andet; forholdet mellem dem kan symboliseres ved forholdet mellem
tallene: 99 er mindre end 100, men 100 indeslutter 99 i sig. De
höjere væsener indeslutte de lavere og desforuden et plus. Således
indgå de enkelte individer som momenter i samfundet som et
mere omfattende individ, og de mindre samfund i de höjere. (”Grundlinier,”
p. 40—41).

Her have vi da den fuldkomne organisme, den ideelle harmoni
skildret, til hvilken vi i det foregående så ofte henvistes af den svenske
filosofis historieskriver. Det er en \textit{mundus intelligibilis,} et
åndigt univers, der her fremstilles for os som den sande virkelighed.
Det er en absolut personlighed --- men netop som absolut personlighed
er den et rige af personlige ånder, et system af absolute væsener.
Den har sit indhold oprindelig i sig; det kan ikke opstå, ikke komme
ude fra, men må være lige så evigt som den selv.

Når det absolute tænkes på denne måde, blive alle de opfattelser
at afvise, der lade dets indhold opstå ved skabelse eller udvikling
og derved endeliggöre det. Når thcologien lader Gud skabe verden,
gör den ham til et endeligt væsen, da skabelsen jo dog må frembringe
en forandring i hans natur, og han altså i forvejen ikke kunde være
absolut fuldkommen. Når den tyske spekulation lader verdensindholdet
opstå ved en proces, der er en evolution af det absolutes eget
% --- p. 52 ---
\setcounter{footnote}{0}%
\opage{52}væsen, bringer den ligeledes endelighed ind i det og indvikler sig derved
i modsigelser. Både skabelses- og produktionstheorien ere endnu befængte
med empirisme, idet de betragte bevægelsen som væsenligt element
i livets og bevidsthedens natur. Det er netop tegn på relativ
mangel eller negation af liv, når levende væsener bevæge sig. Det
i alle henseender levende og selvbevidste væsen kan ikke og behøver
ikke at udvikle og forandre sig, som det endelige må göre. Det er
kun selvstændigheden i bevægelsen, der væsenlig tilhører livet; men
den fuldkomne selvstændighed findes kun hos det absolute. (”Grundlinier”
p. 1—2).\footnote[1]{Polemiken mod den kirkelige theologi og den tyske spekulation findes i de to hvasse stridsskrifter: ”Anmärkningar om Helfveteläran” og ”Den speculativa filosofen J. J. Borelius i Calmar”.}

Dersom man vilde kalde en sådan anskuelse for pantheisme, svarer
Boström dertil, at man ved dette ord dog sædvanlig forstår en
lære, der ikke tænker sig Guddommen som selvstændig og uafhængig
af den sanselige verden, men forestiller sig hin som dennes substans,
og denne som hins accidens, så at bægge gensidig forudsætte hinanden.
”Efter min anskuelse derimod kan den sanselige verden vistnok ikke
existere uden Guddommen, men denne forudsætter og behøver lige så
lidt hin, som f. ex. den virkelige sol behøver den synlige sol. For
mig er næmlig det åndelige det eneste virkelige og oprindelig værende;
og da jeg tænker mig Guddommen som den absolute ånd, må jeg tænke
alt sandt værende i den og ikke sætte den i umiddelbar relation til
noget ydre. Den sanselige verden derimod opfatter jeg som kun udgörende
et fænomen for os af den åndelige og oversanselige, samt
stående omtrent i samme forhold til denne, som f. ex. skumringen
til dagslyset, eller den synlige sol til den virkelige. --- Dersom dette
kan kaldes pantheisme, så er det i hvert tilfælde ikke den sædvanlige”.
(”Satser om lag och lagstiftning”. Andra uppl. Upsala 1871. p. 54).

Der gives altså et andet univers end hint ideelle, uendelig
harmoniske og fuldendte. Men spörge vi, hvor dette sanselige univers,
som er i tid og rum, som forandres, udvikler og opløser sig, da
findes, og hvorledes det opstår, --- da svarer Boström, at de endelige
væsener netop på grund af deres endelighed ikke kunne opfatte
det absolute (idealverdenen) med absolut klarhed og tydelighed,
hvorfor det for dem tager sig ud som noget andet, end
det i sig selv er. Det er kun den uendelige fornuft eller Gud, der
kan opfatte sig selv og derigennem alt andet klart og sandt. Alle
andre væsener må fornemme sig selv og derigennem alt andet mere
eller mindre dunkelt og uklart og for så vidt relativt og ufuldkomment.
Dette har sin grund i, at ethvert andet væsen end det höjeste er endeligt
og derfor efter sit indhold ufuldkomment i sammenligning med
de höjere og det höjeste: det har et ringere mål af opfattelsesævne,
% --- p. 53 ---
\setcounter{footnote}{0}%
\opage{53}på grund af dets plads i væsenernes system. Intet af det evige indhold
går derfor tabt for de endelige væsener, men der kastes en skygge
over det; det fremtræder i uklar, sanselig form. Det sanselige er
kun det fornuftige i dunkel, potentiel, uudviklet (oåtskild) form. Hvert
enkelt endeligt væsen befinder sig på sit ejendommelige stade og har
derfor sit eget individuelle univers, som er mere eller mindre overensstemmende
med de andre væseners, alt efter som der er större eller
mindre slægtskab i deres natur.\footnote[1]{På lignende måde har enhver i ethisk henseende sin egen lov og kan ikke følge nogen anden (”Satser” etc., p. 17---21), eftersom enhver lov, han anerkender, anerkendes af ham, fordi den udtrykker hans væsen, hvad enten dette nu er sanseligt eller fornuftigt.} Der er altså lige så mange verdener,
som der er endelige væsener, og hver verden har sin bestemte
grad af klarhed. Erkendelsens fremskridt består i, at dunkelheden
efterhånden viger, og det evige og uendelige indhold også opfattes
som evigt og uendeligt. Den eneste udvikling, som finder sted, er
denne de endelige væseners overgang fra dunkelhed til klarhed.
Denne udvikling foregår gennem en række af stedse höjere livsformer.
Men selv i den höjeste livsform er sanseligheden --- og altså
dunkelheden --- ikke forsvunden. Den er vel mindre end ved udviklingens
begyndelse, men dersom den fuldstændig ophævedes, vilde
de endelige væsener miste deres selvstændige beståen. Ethvert væsens
udvikling standser, når den for det höjest mulige fuldkommenhed er
nået; og der bliver stedse en forskel mellem systemet af de endelige
væsener i deres höjeste livsformer og systemet af de rent ideelle væsener
eller ideerne, således som de tænkes af Guddommen selv og
ere dens selvbevidste bestemmelser. (”Grundlinier” p. 44 ff, 64 ff).

Man vil let se, at for de endelige væseners vedkommende anerkender
Boström relativitetens lov; for enhver af dem existerer der kun
en relativ virkelighed, og de kunne i al evighed ikke komme ud over
denne relativitet. Derimod antager han, at det forholder sig anderledes
for det absolute væsens vedkommende: det ser tingene, som
de ere i sig selv. Men ligger det ikke ligefrem i selve erkendelsens
begreb, at dette er umuligt? Erkendelse forudsætter en gensidig bestemmelse
af subjekt og objekt; men bestemmelse vil sige relation,
og det, som står i relation til et andet, er forskelligt derfra. Dette
udtaler Boström selv. ”Det er klart”, siger han (”Grundlinier” p. 38),
”at intet kan være, og være noget i og for en anden, når det ikke
først (i begrebet) er, og er noget i og for sig selv; ti det, der ingen
ting var i og for sig selv, vilde ikke kunne tænkes at stå i \textit{forhold til}
noget andet --- hvilket er nødvendigt, når det skal tænkes som en af
dets bestemmelser”. Når altså forhold, relation ikke kan holdes ude,
kan da relativiteten holdes ude? Når erkendelse stedse er bestemt
ved det erkendende væsens natur, så må Guds erkendelse også være
% --- p. 54 ---
\setcounter{footnote}{0}%
\opage{54}bestemt ved hans natur; han må have sin fænomenalverden, ligesom
vi have vor, og han kan --- som Jouffroy meget rigtig har bemærket\footnote[1]{”Dictionnaire des sciences philosophiques” 2:e éd., p. 830.} ---
lige så lidt som vi vide, om hans erkendelse stemmer med den
absolute virkelighed, eller rettere: den ”absolute virkelighed” bliver,
hvor höjt vi end gå til vejrs i de virkelige eller tænkte væseners række,
med hvor ideale attributer vi end udstyre dem, stedse dog kun et
ideal, en grænse, der stadig fjærner sig.

Sagen er, at medens Boström klart har indset relativiteten af begreber
som tid, rum, kraft, udvikling o. s. v., og derfor søger at holde
dem borte fra det absolute, har han ikke set, at erkendelsens og bevidsthedens
begreber ligeledes ere relative, og at de derfor ophæves,
når man forsøger at give dem absolut betydning.

Men en ikke mindre vanskelighed opstår, når vi for et öjeblik
anerkende idealverdenens realitet, --- den samme vanskelighed, som
allerede var til stede for Plato, og som idealismen aldrig har kunnet
overvinde. Hvorfra kommer næmlig den dobbelte (egenlig uendelig
flerdobbelte) opfattelsesmåde af det ene og samme univers? Boström
svarer, at det er endeligheden, der forårsager den dunklere opfattelse.
Og selve endeligheden? Boström svarer (”Anmärkningar” etc., p. 14):
”At spörge efter en årsag til menneskeåndens endelighed vilde omtrent
være lige så urimeligt som at spörge efter en årsag til trianglens tresidighed.
Den vilde jo ikke være en triangel, når den ikke var tresidet;
og ligeså vilde hin ikke være en menneskeånd, når den ikke
var endelig”. Altså må man her ikke længer spörge \textit{hvorfor;} vi stå
ved noget, som ikke længer kan begrundes. Det lavere mål af opfattelsesævne,
som det endelige væsen har, beror på dets ”evige plads
i væsenernes system”. (”Grundlinier” p. 45). Ud over modsætningen
mellem det uendelige og det endelige kunne vi altså aldrig komme:
den er evig.

Og dog er det for at forklare den sanselige verden og den endelige
udvikling, at tanken har konstrueret sig hin ideale verden! Det
er en maxime hos Boström, at det höjere forklarer det lavere (”Grundlinier”
p. 65 f), at det fuldkomne forklarer det ufuldkomne, så at der
ikke kunde gives ufuldkommen erkendelse, når der ikke gaves en
fuldkommen erkendelse. (”Om relig.” etc., p. 47). Men have vi i
virkeligheden vundet nogen forklaring ved at gå tilbage til idealverdenen?
Vi have kun fordoblet problemet, ligesom Plato mente at
forklare de virkelige ting ved at danne en idé for hver af dem.

Måske vil man holde sig til Boström's billede: at sanseverdenen
forudsætter idealverdenen, ligesom den synlige sol forudsætter den
virkelige sol, og dæmringen forudsætter det fulde lys. Men dette
billede forekommer mig at være meget skæbnesvangert for den Boströmske
anskuelse. Ti der vilde ikke være forskel mellem solen som
% --- p. 55 ---
\setcounter{footnote}{0}%
\opage{55}klode og den lille skive, vi se på himlen, når der ikke foruden solkloden
var en anden klode, hvorfra den kunde tage sig ud som en
lille skive; men hvor er et sådant stade muligt, når idealverdenen er
den eneste virkelighed, og der intet er foruden den? Og for at lyset
kan dæmpes, forudsættes der en positiv hindring; men hvor er der
en sådan positiv hindring for det absolute? Boström lader rigtignok
dæmringen være en enhed af lys og mørke; men mørket er ikke andet
end netop lysets negation --- så deraf bliver man ikke klogere. Når
han siger, at ”fænomenalverdenen er et resultat i og for os af to modsatte
faktorers samvirken, en positiv og en negativ” (”Grundlinier” p.
46), så giver dette heller ingen forklaring, da det absolut negative
(= intet) aldeles ikke kan virke, og vi savne oplysning om, hvad det
er for en positiv faktor, der i denne bestemte sammenhæng kan betegnes
som negativ, fordi den virker modsat en anden faktor. Men
en sådan selvstændig positiv faktor vilde stride mod forudsætningen
om idealverdenen som den absolute virkelighed.

Den svenske filosofi fører os skridt for skridt fra empirismens
lave egne opad til spekulationens bjærgtinde. Selv hos Hegel, panlogismens
filosof, finder den endnu ”en god del empirisme”; han skal
endnu ikke have forvundet eftervirkningerne af Locke. Det er i det
hele karakteristisk for den svenske spekulations historiske stilling, at
den stedse opfatter sig selv som modsætning til udviklingsgangen fra
Locke til Hegel, --- som det Elysium, hvor de i denne lange udviklingsperiode
kæmpende tanker kunne hvile ud i uforstyrret harmoni. Spörgsmålet
er, om den ikke ved denne strænge idealisme netop modarbejder
sig selv. Den søger absolute garantier for den menneskelige erkendelse
og handlen, og ingen kan miskende, at der lyder et theoretisk
og praktisk ”excelsior!” igennem den. Men den stiger så höjt op, at
den taber virkeligheden af syne --- virkeligheden, som dog netop
skulde forklares; den fører os stedse höjere op, men kommer derved
til sidst op i den evige snes regioner, hvor jordiske vækster ikke
længer trives. Det er kun i ”erfaringens frugtbare lavland”, vi kunne
bo og bygge; selv om vi imellem stige op på höjderne for at orientere
os og vinde overblik, så opføre vi dog ikke vore huse der oppe.

\begin{center}II.\end{center}

Efter den nu givne historiske  oversigt ville vi, til videre begrundelse
af de lejlighedsvis fremførte kritiske bemærkninger, gå noget
nærmere ind på enkelte spörgsmål, som ere af særegen betydning for
opfattelsen og behandlingsmåden af de filosofiske problemer.

Allerede de første forsøg på at fremstille erkendelsens methode
anerkende, at denne ikke kan være enkelt, men at den består i en
dobbelt bevægelse. Tanken forefinder sit indhold og sig selv som
% --- p. 56 ---
\setcounter{footnote}{0}%
\opage{56}umiddelbar givet, og søger nu at ordne, gruppere, opløse og forbinde
de forskellige dele af dette indhold. Den finder et kaos og vil frembringe
en verden. Den første bevægelse er en analyse, hvorved de
enkelte elementer udskilles af deres første, tilfældige forbindelser;
derpå søger synthesen at bringe nye forbindelser i stand efter elementernes
slægtskab og væsenlige sammenhøren.

Erkendelsen er en kamp med det ubekendte, og ubekendt er oprindelig
alt. Det ubekendte overvindes ved lag for lag at afklædes
sin umiddelbare form og iføres en klædning, som tanken er fortrolig
med. Der er i sig selv intet fast mål sat for analysen; muligheden
af at skride videre frem er bestandig til stede. Der standses, når det
givne er ført tilbage til så simple og gennemskuelige forhold, at det
foreløbig ikke er muligt at nå videre. Synthesen søger da at sammensætte
elementerne til en totalitet. Analysen går fra tingene til
grundene, synthesen fra grundene til tingene. I vekselvirkningen mellem
bægge består videnskabens liv. Goethe har træffende sammenlignet
dem med indånden og udånden. Videnskabens ideal er en fuldstændig
gennemførelse af analysen, hvorved atter en fuldstændig synthese
bliver mulig. Virkelighedens rige og mangfoldige indhold vilde
da være ført tilbage til sit simpleste udtryk, til en formel, hvori dens
hele fylde indesluttedes.

Denne dobbelthed i methoden vil man finde begrundet, ikke blot
ved en undersøgelse af de forskellige videnskabers specielle methoder,
men også ved en betragtning af erkendelsens almindelige væsen og
vilkår. I alt bevidsthedsliv er der en uopløselig forening af enhed
og mangfoldighed, af aktivitet og passivitet. De spredte, isolerede
tilstande blive først bevidste ved at bringes i sammenhæng med hverandre:
ingen fornemmelse er absolut enkelt. Der gives ingen absolute
bevidstheds-atomer. Men enheden er tillige en aktivitet, der lægger
sig for dagen i den uophørlige --- vilkårlige eller uvilkårlige --- sammenligning,
sondren, sammenstillen, opløsning og forbindelse, som foregår
i bevidstheden. Det enkelte indtryk, vi modtage, drages så at
sige ind i en malström straks ved sin indtrædelse i bevidstheden. Den
plads, der anvises det i forestillingskredsen, beror ikke blot på det
selv, men overvejende på dets forhold til andre givne eller reproducerede
forestillinger. Man vil nu let se, at i analysen retter opmærksomheden
sig mod de enkelte bestemmelser og søger at opfatte dem
så skarpt som muligt, hver for sig; i synthesen fremhersker interessen
for atter at restituere enheden og dermed subjektets egen indre aktivitet.
Der bliver altså to poler i erkendelsens verden. Og dette stemmer
fremdeles med erkendelsens betydning for individet. Fra først
af er den intet rent theoretisk anliggende. Det er et spörgsmål på
liv og død, om individet ordner og forbinder sine forestillinger på
rette måde. Så meget derfor end erkendelsen stedse bevarer sin sub%
% --- p. 57 ---
\setcounter{footnote}{0}%
\opage{57}jektive karakter, så må den bestandig stræbe at komme ud over denne,
stræbe at bringe sine forestillingsforbindelser nærmere til den objektive
gyldighed. Erkendelseslivet er kun en enkelt form for den store
kamp for existensen og for den fordring, der stilles til alle organiske
væsener: at de skulle bringe sig i harmoni med deres omverden. En
isoleret subjektivitet vilde være undergangens visse bytte.

Analysen er overalt synthesens grundlag, men det nærmere forhold
mellem de to methoder --- hvorvidt en synthese i det hele er
mulig, og i hvilket omfang --- er forskelligt i de forskellige videnskaber.
Den synthetiske methode kan kun få positiv betydning, når
analysen har ført til simple og ensartede forhold, der i sig selv indeholde
principet for en hel ny og selvstændig udvikling. Dette er
således tilfældet med den analyse, der fører til opstilling af de mathematiske
axiomer. I andre tilfælde bliver synthesen kun en rekonstruktion
ud fra et vist punkt, med stadigt henblik til erfaringen og
med successiv optagen af de ved analysen givne elementer. Virkeligheden
er stedse sammenhæng, konkret totalitet. Men den umiddelbare
sammenhæng står for tanken som tilfældig; det kommer da
an på, om den som frugten af sit arbejde kan sætte en ny totalitet
i stedet for den givne. Ofte går det den som börnene, der rive
deres legetöj i stykker for at få at vide, hvorledes det ser ud indeni,
og nu forgæves prøve på at sætte det sammen igen. Og dog have
börnene fuldstændig ret i, at de ikke kunne lære tingene at kende,
så længe de lade dem blive, som de ere.

Filosofien må være samme lov underlagt, som alle andre videnskabsgrene.
Der gives ingen genveje til erkendelsens mål; vejen er
overalt én og den samme. Den spekulative filosofi har alligevel ofte
taget sit udgangspunkt i almindelige principer og ideer, hvis gyldigheden
ikke betvivlede, og hvis konsekvenser den så forfulgte, så langt
den formåede. Men man vil altid finde, at der er gået et analytisk
arbejde forud --- som jo Sokrates går forud for Plato, Descartes for
Spinoza, Kant for Fichte, Schelling og Hegel. Fejlen er her kun
den, at man for hurtig har anset analysen for fuldendt. På den
anden side må man betragte selve de spekulative systemer som forskellige
forsøg på at undersøge enkelte grundbegreber og prøve deres
rækkevidde. Det enkelte begreb, der göres til systemets princip, er
taget ud af hele den menneskelige idékreds, og på denne måde er
systemernes historie en fremskridende analyse, hvorved anskuelserne
klares og lutres. Dog vægrer i reglen det enkelte system sig ved
således at være led i en fremskridende proces; det mener i sit princip
at være trængt frem til tilværelsens indre og fordrer det anerkendt
som den sande erkendelse. Derved får filosofien så ofte en dogmatisk
karakter, og det vil altid vise sig, at dogmatismen er forbunden
med en undervurdering af den analytiske methode. Som synthetisk
% --- p. 58 ---
\setcounter{footnote}{0}%
\opage{58}videnskab er filosofien system, som analytisk videnskab kritik. På
forholdet mellem det systematiske og det kritiske beror et filosofisk
forsøgs karakter og værdi. Men når bevidstheden engang er vågnet
for betydningen af, at selve principerne undersøges og begrundes,
vil man snart se, at filosofiens videnskabelige karakter beror på den
kritisk-analytiske methode. Den gamle regel de principiis non est
disputandum er, når den skal gælde absolut, aldeles forkastelig. At
man altid i sin tænkning må gå ud fra visse grundsætninger, er sikkert;
men disse grundsætninger selv kan det i mange tilfælde være
fuldt berettiget og nødvendigt at inddrage i diskussionen.

Når vi nærmere ville undersøge, hvorledes filosofien er stillet
med hensyn til forholdet mellem analytisk og synthetisk methode,
må vi sondre mellem de forskellige discipliner, som man i almindelighed
henregner under filosofien. Medens i oldtiden al viden kaldtes
filosofi, har senere efterhånden --- som ved en celledannelse ud af
det strukturløse protoplasma --- den ene specielle videnskab dannet og
udsondret sig efter den anden. På engelsk kaldes fysiken endnu ”natural
philosophy”, som erindring om dens spæde barndom. Det seneste
exempel på en sådan udsondring er den af Adam Smith grundlagte
nationaløkonomi. Nogle af de discipliner, der endnu regnes til filosofien,
have ligesom fagvidenskaberne et aldeles specielt og begrænset
område, og de ville derfor med tiden kunne konstituere sig som selvstændige
videnskaber. Således psykologi, logik og ethik; de synes
næsten at stå på springet. Når således alle enkelte områder få hvert
sin specielle videnskabsgren, sin analyse og synthese, bliver der kun
én opgave tilbage, som tilfalder filosofien i dette ords almindelige betydning
—: den at undersøge, hvilken værdi og hvilken absolut gyldighed
den på alle de specielle områder vundne erkendelse har, på hvilke
forudsætninger den hviler, og hvilken totalopfattelse af virkeligheden
den fører til. Det er altså en analyse i anden potens, der skal foretages;
arbejdet skal ikke göres om, men det skal om muligt underkastes
en reduktion. Og det er dernæst en ny synthese, der skal foretages,
og som skal omfatte alle de specielle syntheser. Det er en opgave,
der med nødvendighed frembyder sig, og som ikke kan skydes
tilbage. Men den er forbunden med ganske særegne vanskeligheder,
som filosoferne ofte have taget det meget let med.

Vi have ovenfor bemærket, at det i synthesen særlig er den subjektive
side i erkendelsen, der spiller en rolle. Den enhed og sammenhæng,
vi søge i fænomenerne, er i sin form bestemt ved vor åndelige
organisation. Synthesen vil derfor let blive mere end en omvendt
repetition af, hvad allerede kom frem under den analytiske undersøgelse;
der bliver et overskud, som skriver sig fra subjektet. Men
hvorledes skal dette plus kontroleres, for at det ikke skal forvanske
det givne? Så længe vi kun bevæge os indenfor en begrænset sfære,
% --- p. 59 ---
\setcounter{footnote}{0}%
\opage{59}er der her ingen uovervindelig vanskelighed. Den sammenhæng, vi
statuere, er da selv led i en större sammenhæng, indenfor hvilken den
står i bestemte forhold; og her vil det da være muligt at prøve, hvorvidt
vor enhedsopfattelse er berettiget eller ikke. Men når synthesen
skal omfatte alt det givne, når vi spörge efter en verdensforklaring,
da er sådan kontrol umulig. Vi have intet, hvormed vi kunne sammenligne
vor totalopfattelse af universet, eller hvortil vi kunne bringe
den i rapport. Den höjeste synthese fører os til et grænsebegreb, til
ideer, der ligge på tænkningens og erkendelsens grænse, ud over hvilken
vi ikke ved noget som helst middel kunne komme. Erkendelsesvirksomheden
er væsenlig betinget ved det bevidste væsens forhold til
den omverden, der gör sig gældende for det; men på det höjdepunkt,
vi her stå ved, falder den objektive side i forholdet bort, og kun den
subjektive side bliver tilbage. Ingen prøvelse ved erfaringen, ingen
verifikation er længer mulig.

Synthesen bliver her nødvendigvis hypothese. Hypothesen har
sin fuldt berettigede plads i videnskaben som forklaringsforsøg, hvor
det foreliggende endnu er utilstrækkeligt til at begrunde en virkelig
erkendelse. Forskningen søger da at samle det givne under et synspunkt,
der så meget som muligt stemmer med kendsgerningerne. Den
fortsætter, kun så at sige med punkterede linjer, de linjer, den er begyndt
at drage indenfor den verificerede erkendelse. Allerede om
nogle af de specielle hypotheser (f. ex. nebularhypothesen, den Darwinske
hypothese) gælder det, at de vanskelig, måske umulig kunne
bevises. Men de forklaringsmuligheder, som der på sådanne specielle
områder vil kunne være tale om, ere dog kun få, og spillerummet vil
stadig indsnævres mere og mere. Derfor kunde man kalde disse hypotheser
reale, i modsætning til dem, der søge at opstille principer
for forklaringen af virkeligheden som helhed. Disse ville stedse beholde
deres ideale og subjektive karakter. Derfor vil tænkerens personlighed
her få en væsenlig indflydelse; den objektive verdensanskuelse
vil modtage et præg af den subjektive livsanskuelse.

Den spekulative filosofi har dog hidtil på ingen måde villet indrömme
dette. Den har stedse i sin höjeste tanke umiddelbart set den
höjeste virkelighed, og den afviser med indignation enhver opfordring
til nærmere bevisførelse; det skal netop være den spekulative erkendelses
særkende, at den har sin sandhed i sig selv, ikke behøver ydre
kriterier på sin gyldighed. Men det sömmer sig ikke for filosofien,
der fremfor alt skulde være selvbevidst erkendelse, og hvis störste
skikkelse gjorde ”kend dig selv!” til sit valgsprog, at dølge den virkelige
sammenhæng for sig selv.

I sin opfattelse af filosofiens methode viser den svenske spekulation
dog större besindighed, end den tyske. Vel anerkender denne
sidste også analysen af det givne som forudsætning for den synthetiske
% --- p. 60 ---
\setcounter{footnote}{0}%
\opage{60}spekulation; men den glemmer snart dette grundlag og griber et princip,
ud fra hvilket den selvstændig kan konstruere sig selv og alt
virkeligt. ”Intet er \textit{a posteriori}, alt er \textit{a priori}, ellers jammerlig halvhed!”
udråbte den energiske Fichte, som var den tyske spekulations
grundlægger. Uden forudsætninger, ganske af egen indre kilde skulde
begreberne udvælde ved en hele tilværelsesindholdet omfattende evolution.
Methode og system skulde være ét, tankens bevægelse ét med
selve indholdets bevægelse.

Men kritiken\footnote[1]{\textit{J. F. Fries} i hans kritik af Schelling (”Reinhold, Fichte und Schelling”; Leipzig 1803). --- \textit{Trendelenburg} i hans kritik af Hegel (”Logische Untersuchungen”, 1840).} påviste, at alle virkelige indholdsbestemmelser
alligevel uden videre (og ofte endda i forvansket form) optoges fra
erfaringen. Det var altså kun illusion, når tanken bildte sig ind, at
den fødte indholdet ud af sig.

Det er netop en af den Boströmske filosofis indvendinger mod
den tyske spekulation, at denne lader det konkrete udvikle sig af det
abstrakte, eller at den forveksler erkendelsesgrund med realgrund.
Det abstrakte, siger prof. Nyblæus i sin kritik af Höijer (som jo i
dette stykke slutter sig til Tyskerne), forudsætter stedse det konkrete.
Kun det absolut konkrete (ɔ: selve det absolute) er forudsætningsløst.
Det konkrete er realgrund for det abstrakte; det abstrakte er for os,
på grund af vor indskrænkning, erkendelsesgrund for det konkrete,
idet vi ved analyse (abstraktion og reflexion) udskille de enkelte bestemmelser
fra det konkrete og undersøge dem for sig. Men af de
således vundne abstrakte principer kunne vi ikke umiddelbar ved synthese
nå det konkrete (de konkretere eller indholdsrigere bestemmelser).
Et begreb kan næmlig aldrig umiddelbar udvikles af et andet,
men [her komme vi til det for den Boströmske erkendelseslære ejendommelige]
kun af de dunklere former for selvbevidsthed, der kaldes
forestilling og følelse. Fremgangen fra videnskabens princip (den
simpleste, mest abstrakte bestemmelse) til det af principet betingede
(”det principierade”) forudsætter derfor, at videnskabens genstand med
hele sit konkrete indhold er given hos tænkeren i form af forestilling
eller følelse, og tænkeren må ved fortsat analyse af dette konkrete
indhold (fra hvilket han allerede havde vundet principet) successivt
udvikle de øvrige, konkretere begreber, der skulle forbindes med principet
og fattes som dets nærmere bestemmelser. Det vil da vise sig,
at der er en bestemt orden, som må følges ved elementernes kombination,
idet den mere abstrakte bestemmelse stedse i begrebet må gå
forud for den mere konkrete som dettes formalgrund. (”Den filosofiska
forskningen”, I, 2 p. 445—447).

Der er i denne fremstilling, i følge sagens natur, to ting at undersøge
nærmere: det givne indhold og dets analyse på den ene side, de
vundne begreber og deres betydning på den anden side.

% --- p. 61 ---
\setcounter{footnote}{0}%
\opage{61}For at en methode som den angivne skal have nogen værdi, er
den første fordring, at indholdet, hvorfra der gås ud, virkelig er givet,
og givet i bestemte former. För enhver uklarhed i denne henseende
er fjærnet, kan analysen ikke med berettigelse afsluttes. Nu er det
umiddelbare indhold, som er givet, når den filosoferende tanke begynder
sit arbejde, af en dobbelt art. Der er for det første den såkaldte
ydre verden, en række billeder og bevægelser, der udfolde sig i tid
og rum, --- og der er for det andet subjektets egne indre tilstande,
dets tanker, følelser og viljesytringer, som afløse hverandre i tiden og
kun symbolsk udformes som billeder i rummet. Men den filosofiske
analyse træffer her ikke på råstof. De enkelte videnskabsgrene analysere
hver sin kreds af den indre og den ydre verdens indhold; dette
arbejdes resultater må den filosofiske analyse anerkende (hvis den da
ikke i spekulativt overmod vil göre det om), og den bliver derfor
--- som vi ovenfor udtrykte os --- analyse i anden potens, når den
forsøger at benytte de specielle analysers resultater til en totalanalyse,
til dannelse af en verdensanskuelse. Men den Boströmske filosofi
synes at mene, at situationen er den samme som i Plato's dage, hvor
filosofen var den eneste forsker, hvor arbejdets deling endnu ikke var
indført på videnskabens område. Den betegner erfaringsvidenskaben
med det hånende ord ”empiri', som for den betegner ”videnskabens”
modstykke. Empirien, siger Boström, har med det sanselige at göre,
spekulationen med det fornuftige, det rationelle. Empirien giver kun
materialerne. En empirisk filosofi har ikke mere ret til filosofiens
navn, end overtroen har til religionens navn. (”Grundlinier” p. 9, 11).
De årsager, som empirien formår at påvise, ere falske, ere kun skinårsager\footnote[1]{”Grundlinier” p. 22 (\dots{}. ”blott falska eller empiriska orsaker --- \textit{causæ efficientes} --- i stället för de sanna eller rationella”), p. 50 (\dots{}. ”skenorsaker, hvilka den empiriska ytligheten anser för verkliga”); kfr. p. 55. --- \textit{Grubbe} stiller ligefrem ”empirisme” i modsætning til ”videnskabelighed”, hvilken falder sammen med spekulationen og den ”högre syftning”. (II p. 323).};
først filosofien finder de sande, rationelle årsager (\textit{causæ finales}). --- Der mangler dog her en væsenlig mellembestemmelse. Polemiken
mod empirien er berettiget, når man ved empiri mener en løs
og tilfældig opsamlen af fakta; allerede Baco polemiserede mod, hvad
han kaldte \textit{experientia vaga}. Men erfaring i betydning af nöjagtig
experimental bestemmelse af lov og sammenhæng er noget helt andet
end hin løse opsamlen, der ganske vist falder udenfor videnskaben.
Den spekulative filosofi gör sig sagen altfor let ved at overse denne
væsenlige forskel. Dersom den vilde stige ned på jorden og erkende
sit fællesskab med den jordiske forsken (det synes ofte, som om den
--- ligesom de til mennesker forvandlede indiske guder --- havde vanskelighed
ved at stå fast på jorden), så vilde den anerkende det som
et stort gode, at erfaringsvidenskaben lægger materialet således til
rette, som den gör. Det givne indhold, filosofien (som spekulation)
% --- p. 62 ---
\setcounter{footnote}{0}%
\opage{62}har med at göre, er altså det i erfaringsvidenskaben givne --- og det
gör en anderledes stærk modstand mod spekulativ behandling, end
\textit{experientia vaga}, hvor der ingen bestemte love er, som kunne brydes.

Det her bemærkede gælder vel i særegen grad om den ydre natur,
men det finder dog også fuld anvendelse på det psykologiske område.
Den empiriske psykologi er endnu så langt tilbage, at man skulde
synes, det måtte være af störste betydning for filosofien at anvende
så mange kræfter som muligt på den for at skaffe sig et solidt udgangspunkt
for sine spekulationer. Hele den nyere filosofis karakter
er betinget ved bevidstheden om nødvendigheden af et sådant psykologisk
grundlag, og denne er også anerkendt i den Boströmske methode,
således som prof. Nyblæus har fremstillet den. Men den vanskelighed,
der så let leder psykologen på fejlspor, næmlig sammenblandingen af
de egenlige psykologiske kendsgerninger med deres forklaring eller
med den værdi, man sætter på dem, --- denne vanskelighed ligger den
Broströmske psykologi i höj grad under for. Den forudsætter, at de
i den umiddelbare følelse ligger et oprindeligt idéindhold skjult, som
det blot gælder at få frem for dagens lys. Den skelner mellem en
lavere og en höjere side i menneskets natur, af hvilke hin skal være
rettet mod det sanselige, denne mod det usanselige. Forskellen imellem
disse to sider er kvalitativ, ɔ: gælder ej blot graden, men art og
indhold. Mennesket har altså ét indhold, én verden i sine lavere følelser,
en anden i sine höjere følelser. (Nyblæus, I, 1 p. 39, 172, 192
o. fl. st.). Men her savner psykologien nöjagtig undersøgelse. Hvad
vil det sige, at en følelse har et indhold? Følelsen er i sig selv kun
følelse, rent subjektiv tilstand; hvad der har et indhold, ere de forestillinger
og tanker, der fremkalde den, eller som den fremkalder.
Men hvorfra komme så igen disse forestillinger, og hvorfra få de deres
jorskellige indhold? Her vil der være plads for mange slags speciellere
undersøgelser, ikke blot psykologiske i snævrere forstand, men
også historiske. Et umiddelbar givet bevidsthedsindhold er aldrig et
sidste, som man uden videre kan lægge til grund; det har sin historiske
oprindelse og udvikling, og først når man kender den, kender
man det. Det er altså meget misligt at drage metafysiske slutninger
fra et sådant bevidsthedsindhold --- og især når æsthetiske og ethiske
værdibestemmelser spille en rolle med. Begreber som ”höjere” og
”lavere” ere netop kun udtryk for den værdi, vedkommende fænomener
kunne have fra et ethisk synspunkt, men bör ikke spille nogen rolle
i psykologien som rent theoretisk disciplin. En sådan inddeling vilde
svare til, om en anthropolog inddelte menneskeracerne eller en botaniker
planterne i skönne og hæslige. Og man kan i hvert tilfælde
ikke fra disse værdibestemmelser slutte, at der er en kvalitativ forskel
mellem fænomenerne: som skönne og hæslige planter kunne være meget
nöje beslægtede, således kunne sjælelige fænomener rent psykologisk
% --- p. 63 ---
\setcounter{footnote}{0}%
\opage{63}stå hinanden meget nær og dog ethisk set henhøre til meget forskellige
klasser.

Dersom analysen skulde have virkelig filosofisk betydning, måtte
erfaringsindholdet først være bredt ud for vore öjne med mere eller
mindre udførlighed, og det måtte være vist, hvorledes de enkelte momenter
kunne udvikles af det. Men vi have allerede set exempler på,
hvorledes begreberne hos Boström undergå en meget betænkelig metamorfose,
når de tages ud af erfaringens kreds. Med den bemærkning,
at bevægelse og forandring ikke ere væsenlige bestemmelser i livets
begreb, mener han at have godtgjort berettigelsen til at anvende dette
som udtryk for det absolute. Et studium af den nyere biologi viser
nu tvært imod, at liv stedse forudsætter en dobbelthed: næmlig foruden
den levende organisme en omverden (le milieu), et system af ydre
betingelser, conditions de l'existence (Cuvier). Og noget lignende gælder
om selve bevidsthedens begreb, således som adskillige filosofer (mest
energisk Fichte) have påvist.\footnote[1]{Boströmianismen opfatter den enkelte personlighed som organ for en större totalitet, en höjere organisme. Men når nu personlighedsbegrebet skulde anvendes på det absolute, så måtte dette jo også være organ for et höjere --- hvilket bliver meningsløst.}

Det turde ved en gennemgående analyse af vort bevidsthedsindhold
vise sig, at der ikke er en eneste bestemmelse, som ikke i sit
væsen er relativ og derfor ikke kan anvendes på det absolute\footnote[2]{Og det hjælper ikke, om man, som Stuart Mill i sin (på dette punkt) uheldige polemik mod Hamilton foreslog, vilde sætte adjektivet ”absolut” i stedet for abstraktionen ”det absolute”. Ti begreber som absolut magt, visdom etc. lide akkurat af samme vanskeligheder, som selve begrebet ”det absolute”; kun skjuler adjektivet dem bedre.}. Det
var indsigten heri, og ikke antagelsen af en ”ting i sig selv” (Nyblæus,
I, 2 p. 253), der førte Kant til at fornægte muligheden af en metafysik
som erkendelse af tingenes væsen. Denne indsigt var endnu levende
hos Fichte og brødes hos ham med den spekulative impetus. Først
Schelling og Hegel satte sig aldeles ud over disse betænkeligheder.
Senere have navnlig Hamilton, Bain, Spencer og Wundt påvist den
gennemgående logiske og psykologiske gyldighed af relativitetens lov.
Kant selv mente vel, at der var visse privilegerede psykologiske
kendsgerninger, der skulde kunne emanciperes og emancipere os fra
denne lov. Pligtfølelsen, det moralske ansvar --- et moment af, hvad
der i Boströmianismen kaldes ”höjere erfaring” --- åbnede for ham
indgangen til det absolutes verden. På lignende måde Hamilton og
flere Ny-Kantianere. Men her gælder det netop, at man ikke må lade
sig lede vild af værdibestemmelserne og ikke uden videre omskrive
dem til theori. Dette er dogmernes frugtbare kilde, en tilflugt for \textit{ignava ratio}. Hvor stor eller hvor liden værdi et faktum har for os,
så kunne vi dog deraf ikke lære det mindste om dets forklaring eller
om, hvorvidt det kan forklares eller ikke forklares af naturlige årsager.
Når den filosofiske diskussion kommer ind på værdibestemmel%
% --- p. 64 ---
\setcounter{footnote}{0}%
\opage{64}sernes område og bruger betegnelserne ”höjere” og ”lavere” som argumenter,
er den inde på et betænkeligt skråplan. Man polemiserer da,
som biskop Mynster hos os, mod Locke som frænde til ”alle slette
og foragtelige filosofemer”, eller man tiltaler sit auditorium, som Victor
Cousin sit i 1815: ”Messieurs, vous aimez ardemment la patrie. Si vous
voulez la sauver, embrassez nos \emph{belles} doctrines”. Tænkningen hører op,
og agitationen begynder. --- Værdibestemmelserne have deres fulde
berettigelse også indenfor filosofien; ethikens og æsthetikens opgave
er det at undersøge og begrunde dem, men de kunne ikke bruges
til at hjælpe på psykologiens og metafysikens (for ikke at tale om
fysikens) vanskeligheder.

Men dersom det nu --- for at komme til det andet hovedpunkt
i den Boströmske methode --- lykkes at udskille grundbegrebernes
række af det empiriske indhold og bringe dem i indre harmoni
med hverandre, hvilken gyldighed kan der da tillægges disse begreber?
--- Når anatomen har udskilt og præpareret forskellige elementer
i organismen, mener han ikke derved at være nået til en ny og höjere
organisme. Ved at betragte de enkelte anatomiske elementer for sig
forstår han bedre deres stilling og rolle indenfor organismen og gör
det lettere for fysiologen at konstruere for sig organismens indre liv.
Men skulde filosofen, der søger at udskille den menneskelige bevidstheds
grundbestemmelser, være anderledes stillet? Han kan ved sin
analyse få dybere indsigt i tankelivets indre sammenhæng, dets udvikling
og dets rækkeævne; men er der noget som helst, der berettiger
ham til at betragte de abstrakte grundbegreber som udtryk for
et realt, höjere og mere omfattende tankeliv? Hvis han gör det, foretager
han her et spring, som han ikke selv mærker eller i hvert fald
ikke forklarer. Det er en mystisk bølge, der griber ham på dette
punkt og fører ham bort med sig. Abstraktion er abstraktion, og kan
aldrig blive andet. Hin hypostaseren af de abstrakte ideer må skyldes
en anden magt end selve den theoretiske erkendelse.

Vi savne på dette punkt en bestemt angivelse af den vej, ad
hvilken erkendelsen formår at nå et positivt begreb om det absolute
som den sande virkelighed. Og dette er så meget mere mærkeligt,
som den Boströmske filosofi ellers lægger megen vægt på bevidsthedens
immanens og så ivrig polemiserer mod den populære realisme,
fordi den går bag ved bevidstheden. Boström giver egenlig kun en
mystisk løsning af denne vanskelighed. Det ligger i bevidsthedens
begreb, siger han, at den kun kan fornemme sig selv og sine egne
fornemmelser. Forudsætningen for, at menneskeånden skal kunne fornemme
Gud, er derfor den, at menneskets egen bevidsthed i en vis
henseende (aliqva ex parte) er identisk med Guds bevidsthed. (”Om
relig.” etc., p. 164). Men hvorfra vides det, at denne forudsætning
finder sted? At min bevidsthed til en vis grad er identisk med et
% --- p. 65 ---
\setcounter{footnote}{0}%
\opage{65}andet væsens bevidsthed, kan jeg kun vide, når jeg kan sammenligne
bægge de to bevidstheder; men så forudsættes, at jeg kender dem
bægge i forvejen. Af denne kreds kunne vi ikke komme ud.

I virkeligheden har al spekulativ filosofi en religiøs karakter. Den
forlanger en opgivelse af reflexionen, for at bevidstheden umiddelbart
kan opfatte sin enhed med det absolute. Derfor hævder den, at der
foruden den diskursive erkendelse gives en höjere erkendelsesart, der
griber det absolute, --- en intuitiv viden, en intellektuel anskuelse, eller
hvorledes man nu betegner den. Og kun derved kan den åbne udsigt
til et höjere virkelighedsindhold, en positiv lære om det absolute. Men
denne fordel købes kun ved opgivelse af, hvad der var vundet ved
kampen mod realismen, idet nu spekulationen selv forsynder sig mod
subjektivitetens lov ved at gå bag ved bevidstheden. Det kommer
næmlig ud på ét, om man statuerer en absolut ånd eller en absolut
materie; man begår i bægge tilfælde samme erkendelsestheoretiske tilsnigelse.
Ligesom Zeus praler af, at han vil drage hele verden op
til sig ved en gylden kæde, knytte denne til en af Olymps tinder og
derved bringe alt i svæven, således har det stedse været spekulationens
opgave at finde et absolut punkt, om hvilket hele tilværelsen
kunde dreje sig. Men vort stade er nu engang ikke det olympiske.
Det punkt, hvortil vi søge at fæste alt, er stedse selv bevægeligt.

Boström indrömmer selv, at bevidstheden, således som vi kende
den, stedse er ufuldkommen og relativ. Den menneskelige stræben
går ud på at nå så stor klarhed og frihed som muligt; med klarheden
og friheden er da forbunden den höjest mulige enhed og harmoni i
alle åndelige funktioner. (”Om relig.”, p. 152). Men ved negation
af vor begrænsning og ufuldkommenhed danne vi begrebet om det
fuldkomne: når det er relationer, der negeres, er negationen en fuldkommenhed.
(”Om relig.”, p. 108). --- Men hvis Boström her ikke
vil sværge til Hegel's fane, er det ham umuligt at nå en ny positivitet
ad negationens vej. Ved negation af det ufuldkomne få vi i virkeligheden
kun det fuldkomne i betydning af det ikke-ufuldkomne. Vi
tænke kun det absolute, for så vidt som vi negere det relative. Men
da der ingen aldeles negativ tænken gives (en sådan vilde være
ét med en ikke-tænken), tager man stedse fantasien til hjælp, når
man vil tænke det absolute. Man forestiller sig mere eller mindre
tydelig et billede, et symbol, som reflexionen öjeblikkelig korrigerer
og opløser, men kun for at det af fantasien kan blive statueret i en
endnu sublimere form, atter opløst --- og således i det uendelige.
Ligesom Antæos stedse hentede ny kraft ved berøring med jorden,
således henter fantasien alle elementer til ideen om det absolute fra
erfaringsverdenen; men tanken fortærer dem ved at hæve dem ud
over den bestemte og begrænsede sfære, indenfor hvilken de have deres
positive betydning.

% --- p. 66 ---
\setcounter{footnote}{0}%
\opage{66}Spekulationens hemmelighed er den, at abstraktionen står i idealisationens
tjæneste. Hver gang vi danne os et ideal, se vi bort fra det
tilfældige og begrænsede, der klæber ved ethvert givet fænomen; vi
holde os til visse egenskaber og karaktermærker, der have vor fulde
tilslutning, og tænke os, hvorledes et væsen vilde være, når det blot
havde disse træk, og havde dem i fuldeste mål og i indbyrdes harmoni.
Denne idealiserende proces kan man tydelig spore i den Boströmske
filosofi. Det åndelige liv, således som erfaringen viser os det, besidder
stedse kun en relativ fuldkommenhed; det kæmper med modsætninger
og modsigelser; det mangler sand enhed og harmoni. Det
höjeste, vi her finde, er en stadig stræben efter fornuft og frihed.
”Menneskeånden besidder et vist mål af liv og selvbevidsthed, ud over
hvilket den ikke i nogen som helst livsform formår at udvikle sig;
hvor höj en grad af fuldkommenhed den så opnår, står der dog stedse
noget ufuldkomnere opfattet tilbage”. (”Om relig.”, p. 162 f.). Erfaringen
viser os altså en skala af udviklingstrin, hvis afslutning er
utænkelig, så længe talen er om endelige væsener. (”Grundlinier” §
56). Men med et rask ryk sætter så fantasien det absolute superlativ
i stedet for disse komparativer; eller rettere: alle komparativerne
ordnes i et system og ses som indbefattede i et evigt superlativ. F.
ex.: mennesket er et fornuftigt væsen, men ikke fornuften selv. Der
er i hans væsen andre elementer end fornuften. Kun Guddommen er
absolut fornuft, fornemmende væsen uden al indskrænkning. (”Satser
om lag och lagstiftning”, p. 22 f.). Vi stå da atter her ved den sætning,
at det ufuldkomne forstås ved det fuldkomne, og vi kunne nu
omskrive den således: Virkeligheden forstås ved idealet.

Der er hos mennesket en trang til at tænke sig formålet for alle
sine bestræbelser færdigt og afsluttet. Han sysler i tanken og længselen
med det, der endnu fattes ham, og som han kæmper for at opnå;
og desuden må han udforme det for sig til en vis grad af bestemthed
for at kunne finde midler og veje til at nå det. Hans nærværende
tilstand er savnets og kampens; men den tilstand, han har for öje, er
det fuldkomment harmoniske og i sig hvilende liv. Og nu vender forholdet
sig således om for ham, at det ser ud, som om formålet var
en allerede tilværende virkelighed, hvorfra han kun skilles ved sin
ufuldkommenhed, så at han kun behøver at tænke denne borte, for at
have ideen om hin. Den ideale tilværelse har for ham i langt dybere
betydning realitet, end den empiriske ufuldkommenhed, som han ikke
værdiger prædikatet af sand virkelighed.

At selve begrebet ”virkelighed” således bliver en værdibestemmelse,
viser, at det er et praktisk motiv, der ligger til grund for spekulationen.
Der kan i de spekulative systemer spores en underström af
følelse og begejstring, som ofte forklarer, hvad systemernes egne begreber
ikke göre forståeligt; og disse systemers blivende værdi beror
% --- p. 67 ---
\setcounter{footnote}{0}%
\opage{67}ikke mindst på, at de ere udtryk for sådanne ubevidste retninger i
følelseslivet. Men heri ligger da også, at sådanne systemer ikke kunne
göre fordring på stræng videnskabelig gyldighed --- end sige, at de
skulde have ret til at betragte sig som udtryk for den sande og höjeste
videnskab.

Prof. Nyblæus, der stærkt fremhæver den svenske spekulations
praktisk-ethiske retning, har påvist, at ideen om den harmoniske totalitet
er dens egenlige grundtanke. Hos Leopold, Thorild, Boëthius
og Höijer se vi den allerede stå i forgrunden, og hos Boström får den
sin fulde udvikling. Det er en ethisk idé af den aller störste betydning.
Det 18de århundredes uhistoriske opfattelse og oppositionelle
retning har bevirket, at man i lang tid miskendte denne idé og troede
at kunne udlede det ethiske af det enkelte individs natur. Men
hvilken side af individets natur man så holder sig til, vil man stedse
føres længere tilbage og opdage, at individet kun forstås som led i
hele slægten, af hvilken det udspringer, hvis arvtager det er, og hvis
opgaver det må arbejde på for at kunne udvikle sit eget væsen. Et
harmonisk samfund af frie personlige væsener kan siges at være det
höjeste ethiske ideal for den menneskelige handlen. Men hvilken ret
have vi til at betragte dette ideal som realt existerende? De er et
forbillede, vi udkaste for vor handlen; vi have det foran os, men det
beror kun på en illusion, når vi tillige sætte det bag os som en absolut
forudsætning. Og ved en sådan hypostasering vil den ethiske interesse
intet vinde, men kun tabe. Den vil intet vinde. Ti det giver
mig ikke nogen dybere indsigt i den ethiske handlens og de ethiske
livsformers væsen, at jeg med Boström tænker mig samfundet som
en höjere, mere omfattende personlighed, hvori jag er moment. Boström
benytter ikke engang selv denne mystiske forudsætning ved
sine undersøgelser. Når han f. ex. vil begrunde sin lære om det
arvelige monarki som den rette statsforfatning, støtter han sig til
grunde, der hentes fra det virkelige samfundslivs krav, idet han navnlig
søger at vise, at regenten må være hævet over alle private og
partikulære interesser, for aldeles uegennyttig at kunne repræsentere
de offenlige og universelle interesser, --- og at statens vel kræver
absolut enhed i styrelsen, hvilket ikke skal kunne opnås i en republik.
Republiken --- mener han --- svarer til den alder hos det enkelte
menneske, hvor lidenskaberne kæmpe med hverandre, medens monarkiet
er udtryk for den fornuftige viljes herredömme. (”Satser om
lag och lagstiftning”, 2. uppl., p. 33 ---35, 63—64). Dersom dette er
rigtigt, så kan det kun bevises med psykologiske og historiske grunde,
og sådanne appelleres der i virkeligheden også bestandig til. På
lignende måde forholder det sig med begrundelsen af strafferetten.
Den Boströmske filosofi godtgör sin humane karakter ved at polemisere
mod den absolute gengældelses-theori, overfor hvilken den hævd%
% --- p. 68 ---
\setcounter{footnote}{0}%
\opage{68}er: at et onde, som straffen jo er, ikke kan være sidste formål, men
kun kan forsvares som nødvendig betingelse for et större gode. Ud
fra denne sætning argumenterer Boström i sin skarpe polemik mod
den theologiske lære om den evige fordömmelse. Og i den ligger
også alt, hvad der behøves til en begrundelse af statens strafferet,
uden at det bliver nødvendigt at opfatte forbrydelsen som krænkelse
af statens mystiske personlighed. I hvert tilfælde er det altid misligt
at bygge på et så hypothetisk og spekulativt grundlag i en så
vigtig sag. Opfattelsen af staten som et harmonisk hele, en total
”personlighed”, angiver kun det mål, der skal nås; midlerne til at nå
dette mål kan kun den fremskridende historiske erfaring lære os, ---
ligesom det jo også er i erfaringen, at de skulle stå deres prøve.\footnote[1]{Smlgn. \textit{Hans Forssell:} ”Studier och kritiker” (Stockholm 1875) p. 174---178.}

Den ethiske interesse vil på den anden side netop svækkes ved,
at dens ideal statueres som en evig væren. Hvorfor skal jeg arbejde,
når alt, hvad jeg tragter efter, allerede er? Det, som er dunkelt,
ufuldkomment og disharmonisk hos mig, er i følge Boströmianismen
kun fenomen, skyldes kun min endelige opfattelse; mit sande væsen
er en evig, harmonisk virkelighed. Men når jeg éngang har indset
dette, ɔ: når jeg er bleven Boströmianer, så kan jeg jo rolig lade hint
fænomen være, hvad det er, for ganske at leve i anskuelsen af mit sande
jeg. Og dersom man vil sige, at modsætningen mellem det klare og det
dunkle, mellem mit væsen og mit fænomen, ikke vil give mig ro dertil,
og at jeg stedse må arbejde theoretisk og praktisk på at ophæve
dunkelheden, så er det ikke godt at indse, hvad hin hypostaseren af
idealet hjælper mig, når det aldeles ikke kan opnås og realiseres af
mig. Selv i sin höjest mulige livsform er det endelige væsen, i følge
Boströmianismen, ved en uoverstigelig kløft skilt fra den absolut
ideale tilværelse. Men hvilken fordel har så egenlig hyperidealisten
fremfor os stakkels empirikere, som vel formå at danne ideale forbilleder,
men ikke kunne tillægge dem anden gyldighed end netop
den at være regler og ledestjærner for vore bestræbelser? Bægge
parter indrömme, at modsætningen mellem ideal og virkelighed er
uovervindelig for os; men hyperidealisten mener desuagtet at kunne
nå en erkendelse af idealet som den eneste sande virkelighed, medens
empirikeren i hin uovervindelighed finder et nyt bevis på relativitetens
lov og et udtryk for, at tilværelsen for os er stræben, bevægelse
og udvikling. Den blivende modsætning mellem ideal og
virkelighed indeholder en brod, som vi bestandig må føle på ny. Det
kan ikke nytte, at Boström trøster os med, at vi dog til sidst nå
målet, for så vidt det kan nås af endelige væsener, og at vi derfor
ikke mere kunne have noget ønske. Psykologien lærer os næmlig,
at så længe vi formå at opfatte forskellen mellem ideal og virkelig%
% --- p. 69 ---
\setcounter{footnote}{0}%
\opage{69}hed, så længe betragte vi også virkeligheden som ufuldkommen. Derfor
står afslutning, fuldendelse, absolut tilfredsstillelse for alle stræbende
naturer som ensbetydende med tilintetgörelse. Lessing's ord er
bekendt: at han vilde foretrække den evige stræben efter sandhed for
besiddelsen af den absolute sandhed, og at den evige salighed vilde
være ét med den evige kedsommelighed. Da Duroc en dag sagde til
Napoleon, at folk anså ham for så ærgerrig, at han gærne, om det
var muligt, vilde indtage Gud Fader's plads, svarede han: ”\textit{Ah! je n'en voudrais \emph{pas; c'est un cul-de-sac}}”. Dette er ikke blot udtryk
for ærgerrig umættelighed, men betegner træffende vor tilværelses
grundlov: ethvert resultat kan for os kun være grundlag for en ny begyndelse.
Hvor langt vi end flytte grænsen, vil det samme gentage sig.

Relativitetens lov lærer os, at vor hele verdensopfattelse er subjektivt
betinget, kun gælder fra vort synspunkt og ikke er den eneste
mulige. Vi prøve så på at hæve os over denne begrænsning, og tro
ved abstraktion fra vor ufuldkommenhed, ved ”negationens negation”
at nå et absolut. Men ethvert nyt begreb, vi ad denne vej tro at
vinde, viser sig stedse at være samme begrænsning og relativitet underlagt.
Det er den gamle historie om nissen, der bestandig flytter med.
Den spekulative filosofi tror, at vor begrænsning er ét, vort sande
væsen et andet, og at man kan udskille hin og holde sig til dette.
Men vi have netop vor styrke i begrænsningen; relativitetens lov er
ikke blot udtryk for vore skranker, men tillige for de positive midler,
hvorover vi råde. Uden skranke og modstand ingen tænken og handlen.
Relativiteten er frugtbar, stræbende og udviklende, hvorimod
absolutismen fører os ind i en theoretisk og praktisk cul-de-sac. Det
absolute er for os kun et \textit{negativt} begreb: vi tænke det, for så vidt
vi tænke vor egen begrænsning. Men som sådant er det af ikke ringe
betydning. Det fører os stedse ud over det umiddelbar givne, og er
derfor princip for stedse fornyet theoretisk og praktisk stræben; vi se
os som ”én blandt mange”, som indtagende et enkelt, begrænset stade
i tilværelsens store mangfoldighed. Gyldigheden af vor theoretiske
indsigt og af vore praktiske idealer kunne vi ikke godtgöre ved at
sammenligne dem med et absolut (hvad enten dette nu er absolut fornuft
eller absolut materie), men kun ved, at de under den stedse fremskridende
erfarings ildprøve vise sig at føre til större sammenhæng og
overensstemmelse mellem vore forestillinger og til större harmoni og
rigere udvikling af det menneskelige liv i det hele. Men en afslutning
på denne proces er for os utænkelig, er ét med den absolute
død, det absolute intet.

Den svenske spekulation er, som vi have set, i sin methode ikke så
dristig som den tyske, stiller sig ikke så store opgaver, som denne gör.
Bægge retninger ere ukritiske ved deres mangelfulde og mislykkede begrundelse
af deres absolute princip; men de skille sig i den videre
% --- p. 70 ---
\setcounter{footnote}{0}%
\opage{70}udførelse, idet Boström mener at kunne hjælpe sig med analysen,
hvorimod Schelling og Hegel fordre en synthetisk konstruktion. Denne
forskel i methoden hænger sammen med en ikke mindre væsenlig
forskel i verdensanskuelsen. I følge Schelling og Hegel er den spekulative
konstruktion ikke blot udtryk for tænkerens subjektive reflexion,
men den falder sammen med det absolutes egen selvudvikling.
Det absolutes væren er en evig proces. Men dette afviser Boström
som empirisme. At vi opfatte tilværelsen som udvikling, skal kun
have sin grund i vor endelighed og begrænsning; vi formå ikke at
fatte det evige idéindhold på en klar og sand måde. Forandringen
og udviklingen angår ikke selve indholdet, men vor opfattelse af det.
Den egenlige udvikling består derfor i vor erkendelses overgang fra
dunkelhed til klarhed, hvorved det evige, absolut oprindelige indhold
mere og mere belyses. Men af denne klaringsproces er det absolute
aldeles uberørt; det er i sig selv stedse fuldendt, medens det sanselige
aldrig i virkeligheden bliver fuldendt. (”Grundlinier” § 1, 6, 35, 38).

Det er ikke blot mod den tyske spekulation, Boströmianismen
her erklærer sig, idet den beskylder den for at göre det absolute endeligt;
det er også mod den kristelige theologi. De kristelige hoveddogmer
(særlig skabelsesdogmet og forsoningsdogmet) hvile på den
forudsætning, at Gud kan virke (producere, skabe) og lide (krænkes og
forsones), og de forsynde sig herved mod det rette og rene begreb
om et absolut væsen. Den spekulative pantheisme er --- som prof.
Nyblæus i tidligere afhandlinger har vist --- kun en ligefrem konsekvens
af de kristelige dogmer. Man kan, siger han, ikke tillægge
Gud producerende eller skabende virksomhed, ikke lade den sanselige
verden være hans umiddelbare værk, uden at hans væsen sammenblandes
med den. Ti hvad der forekommer som resultatet af en virksomhed,
må potentielt være indeholdt i den virkende kraft. Virksomheden
er kun en omsætten fra mulighed til virkelighed af et oprindeligt
indhold. Den, der lader Gud skabe eller producere verden,
gör derfor verden til ét med ham. En sand distinktion mellem det
uendelige og det endelige, det åndelige og det sanselige, bliver kun
mulig, når man antager, at Gud og idealverdenen ere evig fuldendte
og virkelige, og at det kun er i og for vor sanselige opfattelse, der
foregår en udviklingsproces. En uendelig udvikling strider desuden
mod selve udviklingens begreb; ti når målet ligger uendelig langt
borte, kan det aldrig nås, og man savner altså en målestok til at afgöre,
om der virkelig foregår en udvikling fra lavere til höjere trin.
En uendelig kraft vilde lige så lidt kunne tænkes: den vilde straks
være ved målet, dens virkninger vilde falde sammen med den selv\footnote[1]{\textit{Axel Nyblæus:} ”Trenne religionsfilosofiska uppsatser” (Lund 1872) p. 116 f, 126 f, 141, 145. --- ”Om politisk frihet och sjelfstyrelse” (Lund 1869) p. 39, (”Trenne filo%
\opage{71}sofiska uppsatser”, Lund 1878, p. 43). --- ”Den filosofiska forskningen i Sverige”, I, 2 p. 275, 279, 427.}.

% --- p. 71 ---
\setcounter{footnote}{0}%
\opage{71}Det kan sikkert ikke nægtes, at dersom man vil tænke det absolutes
begreb konsekvent og rent, kan der ikke tillægges det udvikling
i tiden, fordi der ikke kan gives noget \textit{plus ultra} for det. Boströmianismen
viser her en betænkelighed, den ikke viser på andre
punkter; ti --- som tidligere bemærket --- de samme vanskeligheder,
som ved dette begreb, stiller der sig for det absolutes filosofi ved begreberne
”liv” og ”selvbevidsthed”. Plotin har set dette: Da tænkning
er en virksomhed, kan den ikke tillægges det höjeste princip. (”Enneades”
III. 9, 3). Er Boström konsekventere end Hegel, så er Plotin
konsekventere end Boström, fordi han har set, at \textit{intet} af vore begreber
slår til overfor fordringen om at tænke det absolute.

Men det er et spörgsmål, om det er berettiget at tillægge Hegel
den lære, at det absolute udvikler sig i tiden. Dets udvikling er i
følge Hegel rent logisk, er kun udtryk for, at det ene moment i dets
begreb med indre nødvendighed fører til de andre, så at på hvilket
punkt man vil begynde, vil man blive ført igennem hele systemet af
dets bestemmelser. Hegel nægter endogså ligefrem en \textit{real} udviklingsproces
i naturen. ”Naturen”, siger han, ”er at betragte som et \textit{system
af trin,} af hvilke det ene nødvendig fremgår af det andet og
er den nærmeste sandhed af det, hvoraf det resulterer, dog ikke således,
at det ene på naturlig måde \textit{blev frembragt} (erzeugt) af \textit{det andet, men
i den indre idé,} der udgör naturens grund”\footnote[1]{\textit{Hegel:} ”Encyclopädie der philosophischen Wissenschaften”, 4te Aufl. (Berlin 1845) p. 204 (§ 249).}. Når spiren udvikler
sig til en plante, udleder Hegel derfor dette af selve spirens logiske
begreb: spiren er den modsigelse, at den kun er i sig selv og dog
ikke skal være det\footnote[2]{\textit{Hegel:} ”Geschichte der Philosophie”, 2te Aufl. (Berlin 1840) I. p. 35.}. Overgangen er altså \textit{egenlig} rent logisk, og
vanskeligheden med tiden kan for så vidt Hegel lige så godt skyde
til side, som Boström gör det, når han tilskriver det absolute tænkning
og selvbevidsthed. Sibbern savnede med rette hos Hegel udviklingen
i betydning af en virkelig \textit{historisk} proces\footnote[3]{\textit{Sibbern:} ”Speculativ Kosmogonie”, (København 1846) p. 7.}.

Selv på det område, hvor Boströmianismen statuerer en udvikling,
har denne ikke nogen historisk karakter. Indholdet er jo evigt og
fuldendt; hvad der går for sig under udviklingen, er kun så at sige
en afdækning eller afskrælning af omgivende hylstre. Boströmianismen
vil ikke indrömme, at der i virkeligheden \textit{sker} noget, ɔ: bliver noget,
som ikke för var. Derfor bebrejder f. ex. prof. Nyblæus Locke, at han
antager en virkelig oprindelse af vore erkendelser i stedet for en
blot overgang fra dunkelhed til klarhed af et oprindeligt, hos os liggende
indhold. (I. 1, p. 3). Hvad der bliver til, skal stedse have været til,
kun fordunklet. Dette peger hen til et udviklingsbegreb, som nu er for%
% --- p. 72 ---
\setcounter{footnote}{0}%
\opage{72}ladt af alle, med undtagelse af spekulative filosofer. Det synes at
være noget, der for mange er forbundet med særegne vanskeligheder,
at indrömme, at der som udviklingens resultat kan fremkomme noget,
der i ingen forstand var til forud. Når barnet om foråret ser bladene
komme frem, tror det, at de have ligget inde i grenene, og at der nu
åbnes for dem. Dets ræsonnement er, at for at de skulle kunne komme
ud, må de først være bragte derind. Udtrykt på mere lærd vis
vil dette sige, at hvad der er aktuelt, allerede må have været til potentielt.
På denne måde får alt en evig præexistens. Således antog
Leibnitz, at fornuften oprindelig var lagt ind i de spirer, der udvikle sig
til fornuftvæsener, så at udvikling (développement) forudsætter indvikling
\textit{(enveloppement}\footnote[1]{”\textit{J'aimerois me passer des miracles dans la génération de l'homme, comme dans celle des autres animaux; et cela se pourra expliquer, en concevant que dans ce grand nombre d'âmes et d'animaux, ou du moins de corps vivans qui sont dans les semences, ces âmes seules qui sont destinées à parvenir un jour à la nature humaine,} \emph{\textit{enveloppent}} \textit{la Raison qui y paraîtra un jour}”. --- ”Théodicée” § 397; kfr. ”Monadologie” § 73---75.}. Det er denne theori, man har kaldet æskehypothesen,
fordi den tænker sig tingene gemte indeni hinanden og udviklingen
bestående i, at man tager dem fra hinanden. Videnskabernes
senere fremskriden har ført til at forlade dette udviklingsbegreb, som
ophæver al kvalitativ forandring. Kant var den første, der opstillede
en anden type for udvikling, idet han i sin ”Allgemeine Naturgeschichte
und Theorie des Himmels” (1755) opstiller den bekendte, endnu gældende
hypothese om solsystemets oprindelse af en kaotisk tågemasse
ved fremskridende artikulation. Kort efter (1759) påviste anatomen
C. F. Wolff, den videnskabelige embryologis stifter, at fosteret opstår
successivt ud fra et ensartet anlæg, og opstillede theorien om epigenese,
om virkelig nydannelse, i modsætning til den såkaldte evolutionstheori,
der lærte en blot afdækning af allerede formede dele. Denne theori
er --- trods Haller og Bonnet --- trængt igennem og har fået sin fulde
bekræftelse\footnote[2]{Smlgn. \textit{Kölliker:} ”Entwicklungsgeschichte des Menschen und der höheren Thiere” (Leipzig 1861) p. 4.}. Den, der vil danne sine begreber efter erfaringen, har
ingen anden udvej end at følge det således indslagne spor. Herbert
Spencer har vist, hvorledes udviklingsbegrebet i sine hovedtræk er
ét og det samme på alle de forskellige områder, og at dets væsenlige
karaktermærker ere overgang fra ensartet til uensartet (differentieret)
tilstand, fra mindre til större sammenhæng, fra ubestemthed til bestemthed\footnote[3]{Se min bog: ”Den engelske Filosofi i vor Tid” (København 1874) p. 136---138, 160---167.}. Den åndelige og den materielle verden, således som vi
kende dem ad erfaringens vej, vise den samme type i deres udviklingsproces.

Overfor den rigdom af træk og begivenheder, udviklingen i erfaringens
verden viser os, tager det sig fattigt ud, når Boströmianismen
ikke har andet at fortælle os, end at det hele er en overgang fra dunkel%
% --- p. 73 ---
\setcounter{footnote}{0}%
\opage{73}hed til klarhed. Det er en barnlig spekulations forsøg på at slippe
fra det paradox --- et paradox, som går gennem hele vor erkendelse
fra først til sidst ---, at noget virkelig bliver til, tanker og erkendelser
så vel som planter, dyr og verdenskloder.

Den tyske spekulation viser sig her at stå erfaringen nærmere,
end den svenske. Den er så levende overbevist om udviklingens betydning,
at den --- ganske vist inkonsekvent --- overfører den på det
absolute. Kristendommens lære om en lidende Gud peger i samme
retning, udtrykker, at det engang er en grundlov, at fuldendelsen
kun nås gennem udvikling, ikke evig er. Fra det praktiske synspunkt
set, er denne overbevisning sikkert af langt större betydning (trods
sin metafysiske inkonsekvens) end Plato's og Boström's rene, af al
vorden, al strid og lidelse uberørte idéverden. Vi ere selv stridende
og lidende væsener, og det nytter os kun lidet at vide, at der er en
sådan evig idealverden, når den stedse er hævet over den bane, på
hvilken vi bevæge os. Det praktiske motiv, hvoraf den svenske spekulation
er udsprungen, har her skudt forbi målet og modarbejder sig
selv. Den verdensanskuelse, der støtter sig til relativitetens grundlov,
formår derimod at holde den idealiserende tendens i den menneskelige
ånd, denne kilde til al poesi og fri spekulation, indenfor de rette
grænser, hvorved den ene bliver i stand til at arbejde i de virkelige
formåls tjæneste.

\end{document}
